
%% LyX 2.0.0 created this file.  For more info, see http://www.lyx.org/.
%% Do not edit unless you really know what you are doing.
\documentclass[oneside,english]{amsart}
\usepackage[T1]{fontenc}
\usepackage[latin9]{inputenc}
\usepackage{amsthm}
\usepackage{amstext}
\usepackage{amssymb}
\usepackage{esint}
\usepackage{graphicx}
\usepackage{color}


\makeatletter
%%%%%%%%%%%%%%%%%%%%%%%%%%%%%% Textclass specific LaTeX commands.
\numberwithin{equation}{section}
\numberwithin{figure}{section}
\theoremstyle{plain}
\newtheorem{thm}{\protect\theoremname}[section]
  \newtheorem{prop}[thm]{\protect\propositionname}
  \theoremstyle{plain}
  \theoremstyle{definition}
  \newtheorem{defn}[thm]{\protect\definitionname}
  \theoremstyle{plain}
  \newtheorem{lem}[thm]{\protect\lemmaname}
  \newtheorem{cor}[thm]{\protect\corollaryname}
  \theoremstyle{plain}
	\newtheorem{rem}[thm]{\protect\remarkname}
  \theoremstyle{plain}
	\theoremstyle{plain}
	\newtheorem{exa}[thm]{\protect\examplename}
\makeatother

\usepackage{babel}
  \providecommand{\definitionname}{Definition}
  \providecommand{\lemmaname}{Lemma}
  \providecommand{\theoremname}{Theorem}
  \providecommand{\corollaryname}{Corollary}
  \providecommand{\remarkname}{Remark}
  \providecommand{\propositionname}{Proposition}
  \providecommand{\examplename}{Example}
	\providecommand{\includegraphicsname}{Fig}
	
	
	\DeclareMathOperator{\Lip}{Lip}
	\DeclareMathOperator{\loc}{loc}
	\DeclareMathOperator{\adj}{adj}
	\DeclareMathOperator{\ess}{ess}
	\DeclareMathOperator*{\esssup}{ess\,sup}
	\DeclareMathOperator{\cp}{cap}
  \DeclareMathOperator{\ACL}{ACL}
	\DeclareMathOperator{\dist}{dist}
  \DeclareMathOperator{\diam}{diam}
  \DeclareMathOperator{\cal}{cal}
  \DeclareMathOperator{\supp}{supp}
  
	
	
\begin{document}

\title[Geometric theory of composition operators on Sobolev spaces]{Geometric theory of composition operators on Sobolev spaces}


\author{Vladimir Gol'dshtein and Alexander Ukhlov} 


\begin{abstract}
In this paper we give the basic concepts of the geometric theory of composition operators on Sobolev spaces. The main objects are topological mappings which generate the bounded embedding operators on Sobolev spaces by the composition rule. This theory is in some sense a "generalization" of the theory of quasiconformal mappings, but the theory of composition operators is oriented to its applications to the Sobolev embedding theorems, the spectral theory of elliptic operators and the continuum mechanics problems.
\end{abstract}
\maketitle
\footnotetext{\textbf{Key words and phrases:} quasiconformal mappings, Sobolev spaces.} 
\footnotetext{\textbf{2010 Mathematics Subject Classification:} 46E35, 30C65.}

\section{Introduction}

This review is devoted to the geometric theory of composition operators on Sobolev spaces. Composition (substitution) operators on Sobolev spaces arise in the Sobolev embedding theory \cite{S41} and in the quasiconformal mapping theory \cite{L71,N60} and form the significant part of the geometric theory of the Sobolev spaces.

The geometric mapping theory of composition operators on Sobolev spaces goes back to the Reshetnyak problem (1968) on geometric characterizations of isomorphisms $\varphi^{\ast}$ of seminormed Sobolev spaces $L^1_n(\Omega)$ and $L^1_n(\widetilde{\Omega})$ generated by the composition rule $\varphi^{\ast}(f)=f\circ\varphi$). Here $\Omega$ and $\widetilde{\Omega}$ are domains in the $n$-dimensional Euclidean space $\mathbb R^n$, $n\geq 2$. In the framework of this problem in \cite{VG75} was proved that these isomorphisms can be generated only by quasiconformal mappings $\varphi:\Omega\to\widetilde{\Omega}$. 

This approach can be considered in the context of the previous knowledge about composition operators on spaces of continuous functions which are Banach algebras (see, \cite{DSh58}). In results by Banach, Stone, Eilenberg, Arens and Kelly, Hewitt, Gel'fand, and Kolmogorov, were obtained characterizations of  various structures of the space of continuous functions $C(S)$, whose isomorphism defines the topological space $S$ up to a topological mapping. On another hand, M. Nakai \cite{N60} and L. Lewis \cite{L71} established an equivalence between isomorphisms of the Royden algebras and corresponding quasiconformal mappings.

Sobolev spaces are not Banach algebras and characteristics of corresponding composition operators are completely different due to the global properties of Sobolev spaces. Consequently, the standard ideas of category theory are also inappropriate in the case of Sobolev spaces $L^1_p(\Omega)$ and $L^1_p(\widetilde{\Omega})$. In \cite{VG76} it was proved that a homeomorphism $\varphi : \Omega\to\widetilde{\Omega}$ generates an isomorphism of Sobolev spaces $L^1_p(\Omega)$ and $L^1_p(\widetilde{\Omega})$, $p>n$, if and only if $\varphi$ is a bi-Lipschitz mapping. This result was extended to the case $n-1<p<n$ in \cite{GR84} and to the case $1\leq p<n$ in \cite{M90} (see, also, \cite{V05}).

Mappings that generate isomorphisms of Besov spaces were considered in \cite{V89}, Nikolskii-Besov spaces and Lizorkin-Triebel spaces were considered in \cite{V90}. In \cite{MS86} the theory of multipliers was applied to the change of variable problem in Sobolev spaces.

The bounded composition operators on Sobolev spaces goes back to the Sobolev embedding theory \cite{S41}. In this paper was proved that smooth bi-Lipschitz homeomorphisms $\varphi: \Omega\to\widetilde{\Omega}$ generates by the composition rule $\varphi^{\ast}(f)=f\circ\varphi$ a bounded operator
$$
\varphi^{\ast}: L^1_p(\widetilde{\Omega})\to L^1_p(\Omega),\,\,\,1\leq p\leq\infty,
$$
and as a consequence, that Sobolev inequalities are preserved under the smooth bi-Lipschitz change of variables.

The composition operators on Sobolev spaces are not necessary isomorphisms even in the case of smooth one-to-one topological mappings.

The next step in the theory of composition operators was given in \cite{GS82} where  extension operators in H\"older cusp domains were constructed using bounded composition operators of $L^1_p(\widetilde{\Omega})$ to $L^1_q(\Omega), q\leq p$ generated by reflections $\varphi:\Omega\to\widetilde{\Omega}$. This work initiated the development of the geometric analysis of Sobolev spaces in non-Lipschitz domains and required the characterization of homeomorphisms $\varphi: \Omega\to\widetilde{\Omega}$ that generate bounded operators
$$
\varphi^{\ast}: L^1_p(\widetilde{\Omega})\to L^1_q(\Omega),\,\,\,1\leq q\leq p\leq\infty.
$$
by the standard composition rule $\varphi^{\ast}(f)=f\circ\varphi$.

The partial solution of this problem, in the case $p=q$ was given in \cite{V88}. The case $q<p$ is more complicated and in this case a solution of this generalized Reshetnyak type problem was given in \cite{U93} in terms of integral characteristics of mappings of finite distortion (the limiting case $p=\infty$ was considered in \cite{GU10-1,GU10-2}) and is based on the countable-additive property of operators norm introduced in \cite{U93}. In the case $p=q=n$ these mappings coincide with quasiconformal mappings. In the general case they are called as weak $(p,q)$-quasiconformal mappings. Note, that the main result of \cite{U93}, in the part of sufficiency of the integral condition, was reproved 20 years later in \cite{Kl12} under some artificial additional conditions. Unfortunately this work \cite{Kl12} does not contains an exact proof due to internal contradictions  between various definitions of Sobolev spaces as Banach spaces used in this paper. Note, that in \cite{Kl12} was not any attempt to prove necessity of the integral conditions.

In \cite{U93} were formulated two main principles of the geometric theory of composition operators on Sobolev spaces:
\begin{itemize}
	\item The localization principle: There exist monotone countable-additive set functions associated with the norms of composition operators generated by mappings restrictions to open subsets.
	\item The composition duality principle: Let a homeomorphism $\varphi: \Omega\to\widetilde{\Omega}$ generate a bounded composition operator
	$$
	\varphi^{\ast}: L^1_p(\widetilde{\Omega})\to L^1_q(\Omega),\,\,\,n-1< q\leq p<\infty.
	$$
	Then the inverse mapping $\varphi^{-1}: \widetilde{\Omega}\to\Omega$ generate a bounded composition operator
	$$
	\left(\varphi^{-1}\right)^{\ast}: L^1_{q'}(\Omega)\to L^1_{p'}(\widetilde{\Omega}),\,\,\,p'=\frac{p}{p-(n-1)}\leq q'=\frac{q}{q-(n-1)}.
	$$
\end{itemize}


In a number of subsequent works \cite{GGR95,GU10-2,U04,VU98,VU02,VU04,VU05} the foundation of the theory of $(p,q)$-quasiconformal mappings were given. The equivalence of the analytic description in the terms of integral characteristics of mappings of finite distortion to the capacitary description and the geometric description, similar the quasiconformal mapping theory is proved. Within the framework of the weak inverse mapping theorem, based on the composition duality property, the regularity of inverse mappings was studied, including the limiting cases $q=n-1$ and $p=\infty$. 

In the case $p=\infty$ homeomorphisms which generate bounded composition operators
$$
	\varphi^{\ast}: L^1_{\infty}(\widetilde{\Omega})\to L^1_q(\Omega),\,\,\,1\leq q\leq \infty,
$$
coincides with the Sobolev homeomorphisms of the class $L^1_q(\Omega;\widetilde{\Omega})$. The regularity of mappings, inverse to Sobolev homeomorphisms $L^1_q(\Omega;\widetilde{\Omega})$, $q\geq n-1$, was intensively studied in the last decades, see, for example \cite{CHM,GU10-1,GU10-2,HK06,HKM06,U04}. 

Unfortunately the proof of the main theorem in \cite{CHM} contains a serious gap. In \cite{CHM} the norm of a weak inverse differential is defined as a composition of a measurable function with the inverse mapping $\varphi^{-1}$. But without the Luzin $N$-property this composition is not necessary a measurable function. It means that the proof in the work \cite{CHM} is correct only under an additional assumption about the Luzin $N$-property of  $\varphi$. Therefore the main theorem in \cite{CHM}  repeats results of \cite{GU10-1,GU10-2} only.

Sobolev homeomorphisms which possess the Luzin $N^{-1}$-property and hence transform by the composition rule measurable functions to measurable functions are called \cite{Z69} as measurable Sobolev homeomorphisms. In \cite{U93,V89,V90} was suggested the approach to measurability of compositions which is based on the non-linear potential theory \cite{HM72}.

The applications of the geometric theory of composition operators on Sobolev spaces to the Sobolev embedding problems are based on the following  "anti-com\-mu\-ta\-tive diagram" which was suggested in \cite{GGu94} 
\[\begin{array}{rcl}
W^{1}_{p}(\widetilde{\Omega}) & \stackrel{\varphi^*}{\longrightarrow} & W^1_q(\Omega) \\[2mm]
\multicolumn{1}{c}{\downarrow} & & \multicolumn{1}{c}{\downarrow} \\[1mm]
L_s(\widetilde{\Omega}) & \stackrel{(\varphi^{-1})^*}{\longleftarrow} & L_r(\Omega)
\end{array}\]
and was used in \cite{GU09} for embeddings of weighted Sobolev spaces. 
In this diagram the operator $\varphi^{\ast}$ defined by the composition rule $\varphi^{\ast}(f)=f\circ\varphi$ is a bounded composition operator on Sobolev spaces induced by a homeomorphism $\varphi$ of $\Omega$ and $\widetilde{\Omega}$ and the operator $(\varphi^{-1})^{\ast}$ defined by the composition rule $(\varphi^{-1})^{\ast}(f)=f\circ\varphi^{-1}$ is a bounded composition operator on Lebesgue spaces. 

This method allows us to obtain Sobolev type embedding theorems in non-Lipschitz domains $\widetilde{\Omega}\subset\mathbb R^n$, that give an approach to the complicated problem of the non-Lipschitz analysis of Sobolev spaces.
By using the Min-Max Principle and the Sobolev embedding theorems, we obtained applications of weak $(p,q)$-quasiconformal mappings to the spectral theory of elliptic operators (see, for example, \cite{GU16,GU17,GPU18_3}). In some cases the composition operators method allows to obtain better estimates than the classical ones by L.~E.~Payne and H.~F.~Weinberger even in convex domains \cite{PW}. In certain sense the method of composition operators on Sobolev spaces represents a remake of the classical method of conformal normalizations \cite{PS51} and gives spectral estimates in non-convex and non-Lipschitz domains even in the case of domains with fractal boundaries. 

The composition operators on Sobolev spaces admit a capacitory description \cite{GGR95,U93,VU98}  and therefore are closely related to so-called $Q$-ho\-me\-o\-mor\-phisms \cite{MRSY}. The  $Q$-homeomorphisms represent a generalization of quasiconformal mappings from the point of view of the geometric function theory. Study of  $Q$-homeomorphisms are based on the capacitary (moduli) distortion of these classes and was under intensive development at last decades (see, for example, \cite{KRSS14,KRSS14M,K86,Sa15,S16}). In the recent work \cite{MU21} it was proved that $Q$-homeomorphisms with integrable function $Q$ coincide with mappings that generate bounded composition operators on Sobolev spaces.
\vskip 0.2cm
 
{\bf Some additional remarks.}
In the paper \cite{IOZ20} (2020)  were considered bi-Sobolev mappings $\varphi:\Omega\to\widetilde{\Omega}$ such that $\varphi\in L^1_n(\Omega)$ and $\varphi^{-1}\in L^1_n(\widetilde{\Omega})$ in connections with the non-linear elasticity problems. Note that these classes were studied in \cite{OS65} as $BL$-mappings and coincide with  classes of weak $(n,1)$-quasiconformal mappings.

In the paper \cite{GPS18} (2018) was stated that the notion of $p$-quasiconformality was introduced in the paper \cite{HK12} (2012) and a version of the composition duality (in the two dimensional case) was introduced in \cite{K12} (2012). The previous papers \cite{GGR95,U93} was not referred in \cite{GPS18}. 

In \cite{MU21} was proved that the class of $Q$-homeomorphisms introduced in \cite{MRSY01} (2001) coincides with the class of weak ($n,n-1$)-quasiconformal mappings. 

\medskip

In this review we give introduction to the geometric theory of composition operators on Sobolev spaces. We consider analytical and capacitory properties of weak $(p,q)$-quasiconformal mappings due to their important applications to the elliptic operators theory. The geometric characterizations of weak $(p,q)$-quasiconformal mappings can be found in \cite{GGR95,VU98}.

Some review on the theory of composition operators on Sobolev spaces can be found in \cite{V12}. Unfortunately, this useful work \cite{V12}  contains some non-correct citations of previous original papers \cite{U93,VU98,VU02}.  


\section{Sobolev Spaces, Capacity and Composition Operators}

\subsection{Sobolev spaces  and Capacity}

The Sobolev space with first weak derivatives $W^1_p(\Omega)$, $1\leq p\leq \infty$ is a space of locally integrable weakly differentiable functions
$f:\Omega\to\mathbb{R}$ equipped with the following norm: 
\[
\|f\mid W^1_p(\Omega)\|=\|f\mid L_p(\Omega)\|+\|\nabla f \mid L_p(\Omega)\|.
\]
It is well known that the Sobolev space $W^1_p(\Omega)$ is a Banach space and coincides with the closure of the space of smooth functions $C^{\infty}(\Omega)$ in the norm of $W^1_p(\Omega)$ \cite{M}. It implies, that functions of the Sobolev space $W^1_p(\Omega)$ are defined (as elements of Banach space) up to a set of $p$-capacity zero, in accordance with the convergence in the topology of  $W^1_p(\Omega)$.

The homogeneous seminormed Sobolev space $L^1_p(\Omega)$, $1\leq p\leq \infty$, is defined as a space
of locally integrable weakly differentiable functions $f:\Omega\to\mathbb{R}$ equipped
with the following seminorm: 
\[
\|f\mid L^1_p(\Omega)\|=\|\nabla f \mid L_p(\Omega)\|.
\]

In the case of Lipschitz domains $\Omega\subset\mathbb R^n$, $n\geq 2$, Sobolev spaces $W^1_p(\Omega)$ and $L^1_p(\Omega)$ coincide because the corresponding Sobolev-Poincar\'{e} inequalities
$$
\inf\limits_{c\in\mathbb R}\|f-c \mid L_p(\Omega)\|\leq B_{p,p}(\Omega)\|f\mid L^1_p(\Omega)\|,
$$
see, for example, \cite{M}. In the general case coincidence of the Sobolev spaces $W^1_p(\Omega)$ and $L^1_p(\Omega)$ is the complicated open problem of the geometric analysis of Sobolev spaces.

As we noted Sobolev spaces are Banach spaces of equivalence classes \cite{M} and it is based on the concept of so-called quasicontinuity which was introduced in 70th by V.~Havin and V.~Maz'ya  \cite{HM72}. This concept is very essential for a correct understanding of the composition operators theory, otherwise it can lead to mistakes like in \cite{HK12} and \cite{K12}.
To clarify the notion of equivalence classes of Banach Sobolev spaces we need to use the non-linear $p$-capacity associated with corresponding Sobolev spaces.

Recall the notion of the $p$-capacity of a set $E\subset \Omega$ \cite{GResh,HKM,M}. 
Suppose $\Omega$ is an open set in $\mathbb R^n$ and  $F\subset\Omega$ is a compact set. The $p$-capacity of $F$ with respect to $\Omega$ is defined as
\begin{equation*}
\cp_p(F;\Omega) =\inf_f\{\|\nabla f|L_p(\Omega)\|^p\},
\end{equation*} 
where the infimum is taken over all functions $f\in C_0(\Omega)\cap L^1_p(\Omega)$ such that $f\geq 1$ on $F$ and which are called admissible functions for the compact set $F\subset\Omega$. 
If 
$U\subset\Omega$ 
is an open set, we define
\begin{equation*}
\cp_{p}(U;\Omega)=\sup_F\{\cp_{p}
(F;\Omega)\,:\,F\subset U,\,\, F\,\,\text{is compact}\}.
\end{equation*}

In the case of an arbitrary set 
$E\subset\Omega$
we define the inner $p$-capacity 
\begin{equation*}
\underline{\cp}_{p}(E;\Omega)=\sup_F\{\cp_{p}(F;\Omega)\, :\,\,F\subset E\subset\Omega,\,\, F\,\,\text{is compact}\},
\end{equation*}
and the outer $p$-capacity 
\begin{equation*}
\overline{\cp}_{p}(E;\Omega)=\inf_U\{\cp_{p}(U;\Omega)\, :\,\,E\subset U\subset\Omega,\,\, U\,\,\text{is open}\}.
\end{equation*}

Then a set $E\subset\Omega$ is called $p$-capacity measurable, if $\underline{\cp}_p(E;\Omega)=\overline{\cp}_p(E;\Omega)$. Let $E\subset\Omega$ be a $p$-capacity measurable set. The value
$$
\cp_p(E;\Omega)=\underline{\cp}_p(E;\Omega)=\overline{\cp}_p(E;\Omega)
$$
is called the $p$-capacity measure of the set $E\subset\Omega$. By \cite{Ch54} analytical sets (in particular Borel sets) are $p$-capacity measurable sets.

The notion of $p$-capacity permits us to refine the notion of Sobolev functions. Let a function $f\in L^1_p(\Omega)$. Then refined function 
$$
\tilde{f}(x)=\lim\limits_{r\to 0}\frac{1}{|B(x,r)|}\int\limits_{B(x,r)} f(y)~dy
$$
is defined quasieverywhere i.~e. up to a set of $p$-capacity zero and it is absolutely continuous on almost all lines \cite{M}. This refined function $\tilde{f}\in L^1_p(\Omega)$ is called a unique quasicontinuous representation ({\it a canonical representation}) of function $f\in L^1_p(\Omega)$. Recall that a function $\tilde{f}$ is termed quasicontinuous if for any $\varepsilon >0$ there is an open  set $U_{\varepsilon}$ such that the $p$-capacity of $U_{\varepsilon}$ is less than $\varepsilon$ and on the set $\Omega\setminus U_{\varepsilon}$ the function  $\tilde{f}$ is continuous (see, for example \cite{HKM,M}). 
In what follows we will use the quasicontinuous (refined) functions only.

Note that the first weak derivatives of
the function $f$ coincide almost everywhere with the usual
partial derivatives (see, e.g., \cite{M} ).

\subsection{The change of variables in the Lebesgue integral}

Let us recall the formula of the change of variables in the Lebesgue integral \cite{Fe69, Ha93}.
Let a mapping $\varphi : \Omega\to \mathbb R^n$ be such that
there exists a collection of closed sets $\{A_k\}_1^{\infty}$, $A_k\subset A_{k+1}\subset \Omega$ for which restrictions $\varphi \vert_{A_k}$ are Lipschitz mapping on the sets $A_k$ and 
$$
\biggl|\Omega\setminus\sum\limits_{k=1}^{\infty}A_k\biggr|=0.
$$
Then there exists a measurable set $S\subset \Omega$, $|S|=0$ such that  the mapping $\varphi:\Omega\setminus S \to \mathbb R^n$ has the Luzin $N$-property and the change of variables formula
\begin{equation}
\label{chvf}
\int\limits_E f\circ\varphi (x) |J(x,\varphi)|~dx=\int\limits_{\mathbb R^n\setminus \varphi(S)} f(y)N_f(E,y)~dy
\end{equation}
holds for every measurable set $E\subset \Omega$ and every nonnegative measurable function $f: \mathbb R^n\to\mathbb R$. Here 
$N_f(y,E)$ is the multiplicity function defined as the number of preimages of $y$ under $f$ in $E$.
If a mapping $\varphi$ possesses the Luzin $N$-property (the image of a set of measure zero has measure zero), then $|\varphi (S)|=0$ and the second integral can be rewritten as the integral on $\mathbb R^n$.

\subsection{Sobolev mappings}

Sufficient conditions on the change of variables formula can be formulated in the terms of Sobolev mappings.
Let $\Omega\subset\mathbb R^n$ be an open set, then a mapping $\varphi:\Omega\to\mathbb R^n$ belongs to Sobolev space $L^1_{p,\loc}(\Omega)$, 
$1\leq p\leq\infty$, if its coordinate functions $\varphi_j$ belong to $L^1_{p,\loc}(\Omega)$, $j=1,\dots,n$.

The Sobolev mappings of the class $W^1_{1,\loc}(\Omega)$ satisfy the conditions of the change of variable formula \cite{Ha93} and so for Sobolev mappings the change of variable formula \eqref{chvf} holds. Note that homeomorphisms of the class $W^1_{n,\loc}(\Omega)$ have the Luzin $N$-property \cite{VGR79}.


In the case of Sobolev mappings the formal Jacobi matrix
$$
D\varphi(x)=\left(\frac{\partial \varphi_i}{\partial x_j}(x)\right),\,\,\, i,j=1,\dots,n,
$$
and its determinant (Jacobian) $J(x,\varphi)=\det D\varphi(x)$ are well defined at
almost all points $x\in \Omega$. The norm $|D\varphi(x)|$ of the matrix
$D\varphi(x)$ is the norm of the corresponding linear operator $D\varphi (x):\mathbb R^n \rightarrow \mathbb R^n$ defined by the matrix $D\varphi(x)$. Recall that the mapping $\varphi$ is called the mapping of finite distortion if $|D\varphi(z)|=0$ for almost all $x\in Z=\{z\in\Omega : J(x,\varphi)=0\}$ \cite{VGR79}.


\section{Composition operators on Sobolev spaces}

Let $\Omega$ and $\widetilde{\Omega}$ be domains in $\mathbb R^n$. Then a homeomorphism $\varphi:\Omega\to\widetilde{\Omega}$ induces a bounded composition
operator 
\[
\varphi^{\ast}:L^1_p(\widetilde{\Omega})\to L^1_q(\Omega),\,\,\,1\leq q\leq p\leq\infty,
\]
by the composition rule $\varphi^{\ast}(f)=f\circ\varphi$, if the composition $\varphi^{\ast}(f)\in L^1_q(\Omega)$
is defined quasi-everywhere (almost everywhere if $q=1$) in $\Omega$ and there exists a constant $K_{p,q}(\varphi;\Omega)<\infty$ such that 
\[
\|\varphi^{\ast}(f)\mid L^1_q(\Omega)\|\leq K_{p,q}(\varphi;\Omega)\|f\mid L^1_p(\widetilde{\Omega})\|
\]
for
any function $f\in L^1_p(\widetilde{\Omega})$ \cite{VU04}.

In the theory of composition operators on Sobolev spaces the significant part consists of set functions defined on open subsets.
Recall that a nonnegative mapping $\Phi$ defined on open subsets of $\Omega$ is called a monotone countably additive set function \cite{RR55,VU04} if

\noindent
1) $\Phi(U_1)\leq \Phi(U_2)$ if $U_1\subset U_2\subset\Omega$;

\noindent
2)  for any collection $U_i \subset U \subset \Omega$, $i=1,2,...$, of mutually disjoint open sets
$$
\sum_{i=1}^{\infty}\Phi(U_i) = \Phi\left(\bigcup_{i=1}^{\infty}U_i\right).
$$

The following lemma gives properties of monotone countably additive set functions defined on open subsets of $\Omega\subset \mathbb R^n$ \cite{RR55,VU04}.

\begin{lem}
\label{lem:AddFun}
Let $\Phi$ be a monotone countably additive set function defined on open subsets of the domain $\Omega\subset \mathbb R^n$. Then

\noindent
(a) at almost all points $x\in \Omega$ there exists a finite derivative
$$
\lim\limits_{r\to 0}\frac{\Phi(B(x,r))}{|B(x,r)|}=\Phi'(x);
$$

\noindent
(b) $\Phi'(x)$ is a measurable function;


\noindent
(c) for every open set $U\subset \Omega$ the inequality
$$
\int\limits_U\Phi'(x)~dx\leq \Phi(U)
$$
holds.
\end{lem}


In the next assertion we formulate the  "localization measure principle" \cite{U93,U04} for composition operators.

\begin{thm}
\label{thm:SobolevAddFun}
 Let a mapping $\varphi : \Omega\to \widetilde{\Omega}$, where $\Omega,\,\widetilde{\Omega}$ are domains of the Euclidean space $\mathbb R^n$,
generates a bounded composition operator
$$
\varphi^{\ast} : L^1_p(\widetilde{\Omega})\to L^1_q(\Omega),\quad 1\leq q<p\leq\infty.
$$
Then
$$
\Phi(\widetilde{A})=\sup\limits_{f\in{L}_{p}^{1}(\widetilde{\Omega})\cap C_0(\widetilde{A})}
\Biggl(
\frac{\bigl\|\varphi^{\ast} f\mid {L}_{q}^{1}(\Omega)\bigr\|}
{\bigl\|f\mid {L}_{p}^{1}(\widetilde{A})\bigr\|}
\Biggr)^{\kappa},
$$
(here number $\kappa$ is defined according to $1/\kappa = 1/q -
1/p$) is a bounded monotone countably additive function defined on
open bounded subsets $\widetilde{A}\subset \widetilde{\Omega}$.
\end{thm}

\begin{proof}
If $\widetilde{A}_1\subset \widetilde{A}_2$ are bounded open  subsets in $\widetilde{\Omega}$, then, extending
the functions of the space $C_0(\widetilde{A}_1)$ by zero on the set $\widetilde{A}_2$, we have the inclusion
$$
\big(L_{p}^{1}(\widetilde{A}_1)\cap C_0(\widetilde{A}_1)\big)\subset
\big(L_{p}^{1}(\widetilde{A}_2)\cap C_0(\widetilde{A}_2)\big).
$$
Hence we obtain the monotonicity of the set function $\Phi$:
\begin{multline}
\Phi(\widetilde{A}_1)=\sup\limits_{f\in
L_{p}^{1}(\widetilde{A}_1)\cap C_0(\widetilde{A}_1)} \Biggl(
\frac{\bigl\|\varphi^{\ast} f\mid {L}_{q}^{1}(\Omega)\bigr\|}
{\bigl\|f\mid
L_{p}^{1}(\widetilde{A}_1)\bigr\|}
\Biggr)^{\kappa}\\
\leq\sup\limits_{f\in
L_{p}^{1}(\widetilde{A}_2)\cap C_0(\widetilde{A}_2)} \Biggl(
\frac{\bigl\|\varphi^{\ast} f\mid {L}_{q}^{1}(\Omega)\bigr\|}
{\bigl\|f\mid
L_{p}^{1}(\widetilde{A}_2)\bigr\|}
\Biggr)^{\kappa}=\Phi(\widetilde{A}_2).
\nonumber
\end{multline}

Now, let $\widetilde{A}_{i}$, $i=1,2,...$, be open disjoint subsets of $\widetilde{\Omega}$ and denote 
$\widetilde{A}_0=\bigcup\limits_{i=1}^{\infty}\widetilde{A}_i$.
For $i=1,2,...$ we consider functions $f_i\in L_{p}^{1}(\widetilde{A}_i)\cap C_0(\widetilde{A}_i)$ to provide simultaneously the following conditions
$$
\bigl\|\varphi^{\ast} f_i\mid{L}_{q}^{1}
(\Omega)\bigr\|\geq \bigl(\Phi(\widetilde{A}_i)\bigl(1-{\frac{\varepsilon}{2^i}}\bigr)\bigr)^{\frac{1}{\kappa}}\bigl\|f_i\mid L_{p}^{1}(\widetilde{A}_i)\bigr\|
$$
and
$$
\bigl\|f_i\mid L_{p}^{1}(\widetilde{A}_i)
\bigr\|^p=\Phi(\widetilde{A}_i)\bigl(1-{\frac{\varepsilon}{2^i}}\bigr),\,\,i=1,2,...,
$$
where $\varepsilon\in(0,1)$ is a fixed number.
Letting $g_N=\sum\limits_{i=1}^{N}f_i$ and applying the H\"older inequality (the case of equality), we obtain
\begin{multline}
\bigl\|\varphi^{\ast} g_N\mid{L}_{q}^{1}
(\Omega)\bigr\| \\ \geq
\biggl(\sum\limits_{i=1}^{N}
\left(\Phi(\widetilde{A}_i)\left(1-{\frac{\varepsilon}{2^i}}\right)\right)^{\frac{q}{\kappa}}
\bigl\|f_i\mid L_{p}^{1}(\widetilde{A}_i)
\bigr\|^q\biggr)^{1/q}\\
=
\biggl(\sum\limits_{i=1}^{N}\Phi(\widetilde{A}_i)
\left(1-{\frac{\varepsilon}{2^i}}\right)\biggr)^{\frac{1}{\kappa}}
\biggl\|g_N\mid L_{p}^{1}
\Bigl(\bigcup\limits_{i=1}^{N}\widetilde{A}_i\Bigr)\biggr\|
\\
\geq
\biggl(\sum\limits_{i=1}^{N}\Phi(\widetilde{A}_i)
-\varepsilon\Phi(\widetilde{A}_0) \biggr)^{\frac{1}{\kappa}}
\biggl\|g_N\mid L_{p}^{1}
\Bigl(\bigcup\limits_{i=1}^{N}\widetilde{A}_i\Bigr)\biggr\|
\nonumber
\end{multline}
since the sets, on which the gradients $\nabla\varphi^{\ast} f_i$
do not vanish, are disjoint. From here it follows that
$$
\Phi(\widetilde{A}_0)^{\frac{1}{\kappa}}\geq\sup\frac
{\bigl\|\varphi^{\ast} g_N\mid{L}_{q}^{1} (\Omega)\bigr\|}
{\biggl\|g_N\mid L_{p}^{1}
\Bigl(\bigcup\limits_{i=1}^{N}\widetilde{A}_i\Bigr)\biggr\|}\geq
\biggl(\sum\limits_{i=1}^{N}\Phi(\widetilde{A}_i)-\varepsilon\Phi(\widetilde{A}_0)
\biggr)^{\frac{1}{\kappa}}
$$
for all $i=1,2,...$, where the least upper bound is taken over all the
above-mentioned functions $g_N\in L_{p}^{1}
\Bigl(\bigcup\limits_{i=1}^{N}\widetilde{A}_i\Bigr)\cap C_0\Bigl(\bigcup\limits_{i=1}^{N}\widetilde{A}_i\Bigr)$. Since both
$N$ and $\varepsilon$ are arbitrary, we have
$$
\sum\limits_{i=1}^{\infty}\Phi(\widetilde{A}_i)
\leq \Phi\Bigl(\bigcup\limits_{i=1}^{\infty}\widetilde{A}_i\Bigr).
$$
The validity of the inverse inequality can be proved straightforward.
\end{proof}


\subsection{Capacitory properties of composition operators}

Let us start with the Luzin type properties of capacities for mappings generating bounded composition operators on Sobolev spaces.

\begin{thm}
\label{theorem:CapacityDescPP}
Let a homeomorphism $\varphi :\Omega\to \widetilde{\Omega}$
generates a bounded composition operator
$$
\varphi^{\ast}: L^1_p(\widetilde{\Omega})\to L^1_p(\Omega),\,\,1\leq p<\infty.
$$
Then the inequality
$$
\cp_{p}^{1/p}(\varphi^{-1}(\widetilde{E});\Omega)
\leq K_{p,p}(\varphi;\Omega)\cp_{p}^{1/p}(\widetilde{E};\widetilde{\Omega})
$$
holds for every Borel set $\widetilde{E}\subset\widetilde{\Omega}$. 
\end{thm}

\begin{proof}

Let $F\subset E=\varphi^{-1}(\widetilde{E})$ be a compact set. Because $\varphi$ is a homeomorphism $\widetilde{F}=\varphi(F)\subset\widetilde{E}$ is also a compact set.  Let $f\in C_0(\widetilde{\Omega})\cap L^1_p(\widetilde{\Omega)}$ be an arbitrary function such that $f\geq 1$ on $\widetilde{F}$. Then the composition $g=\varphi^{\ast}(f)$ belongs to $C_0(\Omega)\cap L^1_p({\Omega)}$, $g\geq 1$ on $F$ and 
$$
\|\varphi^{\ast}(f)\mid L^1_p(\Omega)\|\leq K_{p,p}(\varphi;\Omega)\|f \mid L^1_p(\widetilde{\Omega})\|.
$$
Since the function $g=\varphi^{\ast}(f)\in C_0(\Omega)\cap L^1_p({\Omega)}$ is an admissible function for the compact $F\subset E$, then
$$
\cp_{p}^{1/p}(\varphi^{-1}(\widetilde{F});\Omega)\leq \|\varphi^{\ast}(f)\mid L^1_p(\Omega)\|\leq K_{p,p}(\varphi;\Omega) \|f \mid L^1_p(\widetilde{\Omega})\|.
$$
Taking infimum over all functions $f\in C_0(\widetilde{\Omega})\cap L^1_p(\widetilde{\Omega)}$ such that $f\geq 1$ on $\widetilde{F}$ we have 
$$
\cp_{p}^{1/p}(\varphi^{-1}(\widetilde{F});\Omega)
\leq K_{p,p}(\varphi;\Omega)\cp_{p}^{1/p}(\widetilde{F};\widetilde{\Omega})
$$
for any compact set $\widetilde{F}\subset \widetilde{E}\subset \widetilde{\Omega}$.
Now for the Borel set $\widetilde{E}\subset \widetilde{\Omega}$ we have (by the definition of the $p$-capacity of Borel sets)
$$
\cp_{p}^{1/p}(\widetilde{F};\widetilde{\Omega})\leq \underline{\cp}_{p}^{1/p}(\widetilde{E};\widetilde{\Omega})=\cp_{p}^{1/p}(\widetilde{E};\widetilde{\Omega}).
$$
Hence 
$$
\cp_{p}^{1/p}(\varphi^{-1}(\widetilde{F});\Omega)\leq K_{p,p}(\varphi;\Omega)\cp_{p}^{1/p}(\widetilde{E};\widetilde{\Omega}).
$$
Since $F=\varphi^{-1}(\widetilde{F})$ is an arbitrary compact set, $F\subset E$, $E$ is a Borel set as a preimage of the Borel set $\widetilde{E}$ under the homeomorphism $\varphi$, then
\begin{multline*}
\cp_{p}^{1/p}(\varphi^{-1}(\widetilde{E});\Omega)=\underline{\cp}_{p}^{1/p}(\varphi^{-1}(\widetilde{E});\Omega)=
\sup\limits_{F\subset E}\cp_{p}^{1/p}(F;\Omega)\\
\leq K_{p,p}(\varphi;\Omega)\cp_{p}^{1/p}(\widetilde{E};\widetilde{\Omega}).
\end{multline*}
\end{proof}

\begin{thm}
\label{theorem:CapacityDescPQ}
Let a homeomorphism $\varphi :\Omega\to \widetilde{\Omega}$
generates a bounded composition operator
$$
\varphi^{\ast}: L^1_p(\widetilde{\Omega})\to L^1_q(\Omega),\,\,1\leq q< p<\infty.
$$
Then the inequality
$$
\cp_{q}^{1/q}(\varphi^{-1}(\widetilde{E});\Omega)
\leq K_{p,q}(\varphi;\Omega)\cp_{p}^{1/p}(\widetilde{E};\widetilde{\Omega})
$$
holds for every Borel set $\widetilde{E}\subset\widetilde{\Omega}$. 
\end{thm}

\begin{proof}

Let $F\subset E=\varphi^{-1}(\widetilde{E})$ be a compact set. Because $\varphi$ is a homeomorphism $\widetilde{F}=\varphi(F)\subset\widetilde{E}$ is also a compact set.  Let $f\in C_0(\widetilde{\Omega})\cap L^1_p(\widetilde{\Omega)}$ be an arbitrary function such that $f\geq 1$ on $\widetilde{F}$. Then the composition $g=\varphi^{\ast}(f)$ belongs to $C_0(\Omega)\cap L^1_q({\Omega)}$, $g\geq 1$ on $F$ and 
$$
\|\varphi^{\ast}(f)\mid L^1_q(\Omega)\|\leq K_{p,q}(\varphi;\Omega) \|f \mid L^1_p(\widetilde{\Omega})\|.
$$
Since the function $g=\varphi^{\ast}(f)\in C_0(\Omega)\cap L^1_q({\Omega)}$ is an admissible function for the compact $F\subset E$, then
$$
\cp_{q}^{1/q}(\varphi^{-1}(\widetilde{F});\Omega)\leq \|\varphi^{\ast}(f)\mid L^1_q(\Omega)\|\leq K_{p,q}(\varphi;\Omega) \|f \mid L^1_p(\widetilde{\Omega})\|.
$$
Taking infimum over all functions $f\in C_0(\widetilde{\Omega})\cap L^1_p(\widetilde{\Omega)}$ such that $f\geq 1$ on $\widetilde{F}$ we have 
$$
\cp_{q}^{1/q}(\varphi^{-1}(\widetilde{F});\Omega)
\leq K_{p,q}(\varphi;\Omega)\cp_{p}^{1/p}(\widetilde{F};\widetilde{\Omega})
$$
for any compact set $\widetilde{F}\subset \widetilde{E}\subset \widetilde{\Omega}$.

Now for the Borel set $\widetilde{E}\subset \widetilde{\Omega}$ we have (by the definition of the $p$-capacity of Borel sets)
$$
\cp_{p}^{1/p}(\widetilde{F};\widetilde{\Omega})\leq \underline{\cp}_{p}^{1/p}(\widetilde{E};\widetilde{\Omega})=\cp_{p}^{1/p}(\widetilde{E};\widetilde{\Omega}).
$$
Hence 
$$
\cp_{q}^{1/q}(\varphi^{-1}(\widetilde{F});\Omega)\leq K_{p,q}(\varphi;\Omega)\cp_{p}^{1/p}(\widetilde{E};\widetilde{\Omega}).
$$
Since $F=\varphi^{-1}(\widetilde{F})$ is an arbitrary compact set, $F\subset E$, $E$ is a Borel set as a preimage of the Borel set $\widetilde{E}$ under the homeomorphism $\varphi$, then
\begin{multline*}
\cp_{q}^{1/q}(\varphi^{-1}(\widetilde{E});\Omega)=\underline{\cp}_{q}^{1/q}(\varphi^{-1}(\widetilde{E});\Omega)=
\sup\limits_{F\subset E}\cp_{q}^{1/q}(F;\Omega)\\
\leq K_{p,q}(\varphi;\Omega)\cp_{p}^{1/p}(\widetilde{E};\widetilde{\Omega}).
\end{multline*}

\end{proof}



\begin{cor}
\label{cor:Capacity}
Let a homeomorphism $\varphi :\Omega\to \widetilde{\Omega}$
generates a bounded composition operator
$$
\varphi^{\ast}: L^1_p(\widetilde{\Omega})\to L^1_q(\Omega), \,\,1<q\leq p<\infty.
$$
Then preimage of a set of $p$-capacity zero has $q$-capacity zero.
\end{cor}

\begin{proof}
The proof follows from Theorem~\ref{theorem:CapacityDescPP} and Theorem~\ref{theorem:CapacityDescPQ}.

\end{proof}


The characterization of mappings generate bounded composition operators on Sobolev spaces in capacitary terms is given in the terms of the variational $p$-capacity \cite{GResh}. The condenser in the domain $\Omega\subset \mathbb R^n$ is the pair $(F_0,F_1)$ of connected closed relatively to $\Omega$ sets $F_0,F_1\subset \Omega$. A continuous function $u\in L_p^1(\Omega)$ is called an admissible function for the condenser $(F_0,F_1)$,
if the set $F_i\cap \Omega$ is contained in some connected component of the set $\operatorname{Int}\{x\vert u(x)=i\}$,\ $i=0,1$. We call $p$-capacity of the condenser $(F_0,F_1)$ relatively to domain $\Omega$
the value
$$
{{\cp}}_p(F_0,F_1;\Omega)=\inf\|u\vert L_p^1(\Omega)\|^p,
$$
where the greatest lower bond is taken over all admissible for the condenser $(F_0,F_1)\subset\Omega$ functions. If the condenser have no admissible functions we put the capacity is equal to infinity. 

 Let ${{\cp}}_p(F_0,F_1;\Omega)<\infty$.
Suppose that a function $v$ belongs to $L_p^1(\Omega)$ is admissible for the condenser $(F_0,F_1)$. Then $v$ is called {\it an extremal function for the condenser $(F_0,F_1)$} \cite{VG77}, if
$$
\int\limits_{\Omega\backslash(F_0\cup F_1)}|\nabla v|^p\,dx=
{{\cp}}_p(F_0,F_1;\Omega).
$$

Note that for any $1<p<\infty$ and any condenser $(F_0,F_1)$ with ${{\cp}}_p(F_0,F_1;\Omega)<\infty$ the extremal function exists and unique.

The set of extremal functions for $p$-capacity of every possible pairs of $n$-dimensional connected compacts 
$F_0,\;F_1\subset \Omega$, having smooth boundaries, we denote by the symbol $E_p(\Omega)$.

\begin{thm}
\label{thm:ApproxExrem}
\cite{VG77}
Let $1<p<\infty$. Then there exists a countable collection of functions $v_i\in E_p(\Omega)$, $i\in \mathbb N$,
such that for every function $u\in L_{p}^{1}(\Omega)$ and for any $\varepsilon>0$
there exists a presentation of $u$ in the form
$u=c_0+\sum\limits_{i=1}^{\infty}c_iv_i$,
for which the inequalities 
$$
\|u\mid L_{p}^{1}(\Omega)\|^p\leq \sum\limits_{i=1}^{\infty}
\|c_iv_i\mid L_{p}^{1}(\Omega)\|^p\leq
        \|u\mid L_{p}^{1}(\Omega)\|^p+\varepsilon
$$ 
hold.
\end{thm}

The following assertion was proved, but not formulated in  \cite{VG77}.

\begin{lem}
\label{lem:extremal} Let $1<q\leq p<\infty$. Suppose that $\varphi :\Omega\to \widetilde{\Omega}$ is a homeomorphism of domains $\Omega,\widetilde{\Omega}\subset\mathbb R^n$, such that for a pair of $n$-dimensional connected compacts $\widetilde{F}_0,\widetilde{F}_1\subset \widetilde{\Omega}$ with smooth boundaries the inequality
$$
\left(\int\limits_{\Omega}|\nabla v|^q~dx\right)^{\frac{1}{q}}\leq K_{p,q}(\widetilde{F}_0,\widetilde{F}_1) \left(\int\limits_{\widetilde{\Omega}}|\nabla u|^p~dx\right)^{\frac{1}{p}}
$$
holds for extremal functions $v$ and $u$ of pairs $(\varphi^{-1}(\widetilde{F}_0),\varphi^{-1}(\widetilde{F}_1))$ and $(\widetilde{F}_0,\widetilde{F}_1)$ correspondly, where $K_{p,q}(\widetilde{F}_0,\widetilde{F}_1)$ depends on $\widetilde{F}_0,\widetilde{F}_1\subset\mathbb R^n$. Then
$$
\|\varphi^{\ast}(u)\mid L^1_q(\Omega)\|\leq K_{p,q}(\widetilde{F}_0,\widetilde{F}_1)\|u \mid L^1_p(\widetilde{\Omega})\|.
$$
\end{lem}

\begin{rem}
The proof of this assertion in \cite{VG77} is based on the monotonicity of the capacity and the lower semi-continuity of Sobolev norms.
\end{rem}

In the next assertion gives the capacitory description of composition operators on Sobolev spaces $L^1_p(\widetilde{\Omega})$ and $L^1_p(\Omega)$. In the proof we follows \cite{U93}.

\begin{thm}
\label{thm:CapacityDescPP}
Let $\Omega$ and $\widetilde{\Omega}$ be domains in $\mathbb R^n$. Then a homeomorphism $\varphi :\Omega\to \widetilde{\Omega}$
generates a bounded composition operator
$$
\varphi^{\ast}: L^1_p(\widetilde{\Omega})\to L^1_p(\Omega),\,\,1<p<\infty,
$$
if and only if there exists a constant  $K_{p,p}(\varphi;\Omega)$ such that for every condenser 
$(F_0,F_1)\subset \Omega$
the inequality
$$
\cp_{p}^{1/p}(\varphi^{-1}(F_0),\varphi^{-1}(F_1);\Omega) \leq K_{p,p}(\varphi;\Omega)\cp_{p}^{1/p}(F_0,F_1;\widetilde{\Omega})
$$
holds 
\end{thm}

\begin{proof}
{\it Necessity.}
Let $u$ be an admissible function for the condenser $(F_0,F_1)\subset\widetilde{\Omega}$,
then $u\circ\varphi$ be an admissible function for the condenser
$(\varphi^{-1}(F_0),\varphi^{-1}(F_1))\subset\Omega$.
Since $\varphi$
generates a bounded composition operator
$$
\varphi^{\ast}: L^1_p(\widetilde{\Omega})\to L^1_p(\Omega)
$$
then, by the definition of the $p$-capacity, we have
$$
{\cp}_p^{1/p}(\varphi^{-1}(F_0),\varphi^{-1}(F_1);\Omega)\leq\|\varphi^{\ast}(u)\vert L_q^1(\Omega)\|\leq K_{p,p}(\varphi;\Omega)\|u\,|\,L_p^1(\widetilde{\Omega})\|.
$$
Since $u$ is an arbitrary admissible function, we obtain
$$
\cp_{p}^{1/p}(\varphi^{-1}(F_0),\varphi^{-1}(F_1);\Omega)
\leq K_{p,p}(\varphi;\Omega)\cp_{p}^{1/p}(F_0,F_1;\widetilde{\Omega}).
$$

{\it Sufficiency.}
Let a function $u\in L^1_p(\widetilde{\Omega})$. Then by Theorem~\ref{thm:ApproxExrem} 
$$
u=c_0+\sum\limits_{i=1}^{\infty}c_i v_i,\,\,v_i\in E_p(\widetilde{\Omega}),
$$
and
$$
\|u\mid L_{p}^{1}(\widetilde{\Omega})\|^p\leq \sum\limits_{i=1}^{\infty}\|c_iv_i\mid L_{p}^{1}(\widetilde{\Omega})\|^p\leq
\|u\mid L_{p}^{1}(\widetilde{\Omega})\|^p+\varepsilon
$$ 
for arbitrary $\varepsilon>0$.

The series $c_0+\sum\limits_{i=1}^{\infty}c_i v_i$ converges to the function $u\in L^1_p(\widetilde{\Omega})$ $p$-quasieverywhere (up to a set of $p$-capacity zero) in $\widetilde{\Omega}$. Hence, by Corollary~\ref{cor:Capacity} we have that the series $c_0+\sum\limits_{i=1}^{\infty}c_i (v_i\circ\varphi)$ converges to the function $u\circ\varphi=\varphi^{\ast}(u)$ $p$-quasieverywhere  (up to a set of $p$-capacity zero) in $\Omega$.

By Lemma~\ref{lem:extremal} we obtain
\begin{multline}
\|(c_0+\sum\limits_{i=1}^\infty c_i \varphi^{\ast}(v_i))\vert L_p^1(\Omega)\|^p=
\sum\limits_{i=1}^\infty c_i^p \| \varphi^{\ast}(v_i)\vert L_p^1(\Omega)\|^p\\
\le \sum\limits_{i=1}^\infty K_{p,p}^p(\widetilde{F}_0^i,\widetilde{F}_1^i)c_i^p \| v_i \vert L_p^1(\widetilde{\Omega})\|^p
\leq K_{p,p}^p(\varphi;\Omega) \sum\limits_{i=1}^\infty \|c_i v_i \vert L_p^1(\widetilde{\Omega})\|^p\\
=K_{p,p}^p(\varphi;\Omega) \|u \vert L_p^1(\widetilde{\Omega})\|^p+\varepsilon K_{p,p}^p(\varphi;\Omega).
\nonumber
\end{multline}
Since $\varepsilon>0$ is arbitrary, we obtain,
that for every function $u\in L_p^1(\widetilde{\Omega})$
$$
\|\varphi ^{\ast}(u)\vert L_p^1(\Omega)\|\le K_{p,p}(\varphi;\Omega)\|u\vert L_p^1(\widetilde{\Omega})\|.
$$
\end{proof}

The case $1<q<p<\infty$ was given in \cite{U93}.

\begin{thm}
\label{thm:CapacityDescPQ}
Let $\Omega$ and $\widetilde{\Omega}$ be domains in $\mathbb R^n$. Then a homeomorphism $\varphi :\Omega\to \widetilde{\Omega}$
generates a bounded composition operator
$$
\varphi^{\ast}: L^1_p(\widetilde{\Omega})\to L^1_q(\Omega),\,\,1<q<p<\infty,
$$
if and only if
there exists a bounded monotone countable-additive set function
$\Phi$ defined on open subsets of $\Omega$
such that for every condenser 
$(F_0,F_1)\subset \Omega$
the inequality
$$
\cp_{q}^{1/q}(\varphi^{-1}(F_0),\varphi^{-1}(F_1);\Omega)
\leq\Phi(\widetilde{\Omega}\setminus(F_0\cup F_1))^{\frac{p-q}{pq}}
\cp_{p}^{1/p}(F_0,F_1;\widetilde{\Omega})
$$
holds 
\end{thm}

\begin{proof}
{\it Necessity.}
Let $u$ be an admissible function for the condenser $(F_0,F_1)\subset\widetilde{\Omega}$,
then $u\circ\varphi$ be an admissible function for the condenser
$(\varphi^{-1}(F_0),\varphi^{-1}(F_1))\subset\Omega$.
Since $\varphi$
generates a bounded composition operator
$$
\varphi^{\ast}: L^1_p(\widetilde{\Omega})\to L^1_q(\Omega)
$$
then, by the definition of the $p$-capacity, we have
\begin{multline}
{\cp}_q^{1/q}(\varphi^{-1}(F_0),\varphi^{-1}(F_1);\Omega)\leq
\|\varphi^{\ast}(u)\vert L_q^1(\Omega)\|\\
\le
\Phi(\widetilde{\Omega}\setminus(F_0\cup F_1))^{(p-q)/pq}\|u\,|\,L_p^1(\widetilde{\Omega})\|.
\nonumber
\end{multline}
Since $u$ is arbitrary admissible function, we have
$$
{\cp}_q^{1/q}(\varphi^{-1}(F_0),\varphi^{-1}(F_1);\Omega)
\le
\Phi(\widetilde{\Omega}\setminus(F_0\cup F_1))^{(p-q)/pq}{\cp}_p^{1/p}(F_0,F_1;\widetilde{\Omega}).
$$

{\it Sufficiency.}
Let a function $u\in L^1_p(\widetilde{\Omega})$. Then by Theorem~\ref{thm:ApproxExrem} 
$$
u=c_0+\sum\limits_{i=1}^{\infty}c_i v_i,\,\,v_i\in E_p(\widetilde{\Omega}),
$$
and
$$
\|u\mid L_{p}^{1}(\widetilde{\Omega})\|^p\leq \sum\limits_{i=1}^{\infty}\|c_iv_i\mid L_{p}^{1}(\widetilde{\Omega})\|^p\leq
\|u\mid L_{p}^{1}(\widetilde{\Omega})\|^p+\varepsilon
$$ 
for arbitrary $\varepsilon>0$.

The series $c_0+\sum\limits_{i=1}^{\infty}c_i v_i$ converges to the function $u\in L^1_p(\widetilde{\Omega})$ $p$-quasieverywhere (up to a set of $p$-capacity zero) in $\widetilde{\Omega}$. Hence, by Corollary~\ref{cor:Capacity} we have that the series $c_0+\sum\limits_{i=1}^{\infty}c_i (v_i\circ\varphi)$ converges to the function $u\circ\varphi=\varphi^{\ast}(u)$ $q$-quasieverywhere  (up to a set of $q$-capacity zero) in $\Omega$.

By Lemma~\ref{lem:extremal} we obtain
\begin{multline}
\|(c_0+\sum\limits_{i=1}^\infty c_i \varphi^{\ast}(v_i))\vert L_q^1(\Omega)\|^q=
\sum\limits_{i=1}^\infty c_i^q \| \varphi^{\ast}(v_i)\vert L_q^1(\Omega)\|^q\\
\le \sum\limits_{i=1}^\infty K_{p,q}^q(\widetilde{F}_0^i,\widetilde{F}_1^i)c_i^q \| v_i \vert L_p^1(\widetilde{\Omega})\|^q\leq 
\sum\limits_{i=1}^\infty \Phi(\widetilde{\Omega}\setminus(F_0^i\cup F_1^i))^{\frac{p-q}{p}}c_i^q \| v_i \vert L_p^1(\widetilde{\Omega})\|^q\\
\leq\left(\sum\limits_{i=1}^\infty \Phi(\widetilde{\Omega}\setminus(F_0^i\cup F_1^i))\right)^{\frac{p-q}{p}} 
\left(\sum\limits_{i=1}^\infty c_i^p\| v_i \vert L_p^1(\widetilde{\Omega})\|^p\right)^{\frac{q}{p}}\\
\leq K_{p,q}^{q}(\varphi;\Omega) \left(\sum\limits_{i=1}^\infty \|c_i v_i \vert L_p^1(\widetilde{\Omega})\|^p\right)^{\frac{q}{p}}\\
= K_{p,q}^{q}(\varphi;\Omega) \left(\|u\vert L_p^1(\widetilde{\Omega})\|^p+\varepsilon\right)^{\frac{q}{p}}.
\nonumber
\end{multline}
Since $\varepsilon>0$ is arbitrary, we obtain,
that for every function $u\in L_p^1(\widetilde{\Omega})$
$$
\|\varphi ^{\ast}(u)\vert L_q^1(\Omega)\|\le K_{p,q}(\varphi;\Omega)\|u\vert L_p^1(\widetilde{\Omega})\|.
$$
\end{proof}



\subsection{The analytical description of composition operators}

We begin with an analytical description of homeomorphisms generating bounded composition operators on Sobolev spaces $L^1_p(\widetilde{\Omega})$ and $L^1_p(\Omega)$.

\begin{thm}
\label{thm:AnalyticPP}
Let $\Omega$ and $\widetilde{\Omega}$ be domains in $\mathbb R^n$. Then a homeomorphism $\varphi :\Omega\to \widetilde{\Omega}$ generates a bounded composition operator
$$
\varphi^{\ast} : L^1_p(\widetilde{\Omega})\to L^1_p(\Omega), \,\,\,1\leq p<\infty,
$$
if and only if $\varphi\in L^1_{p,\loc}(\Omega)$ has finite distortion and
$$
\ess\sup\limits_{x\in\Omega}\left(\frac{|D\varphi(x)|^p}{|J(x,\varphi)|}\right)^{\frac{1}{p}}=K_{p,p}(\varphi;\Omega)<\infty.
$$
The norm of the operator $\|\varphi^{\ast}\|\leq K_{p,p}(\varphi;\Omega)$.
\end{thm}

\begin{proof}
{\it Necessity.} 
Fix a cut function $\eta\in C_{0}^{\infty}(\mathbb R^n)$, which is equal to one on the ball $B(0,1)$ and is equal to zero outside of the ball $B(0,2)$.
Now we consider the test functions
$$
f_j(y)=(y_j-y_{j0})\eta\left(\frac{y-y_0}{r}\right),\,\,\, j=1,...,n,
$$
where $y_j$ is the $j$-th coordinate function. We see that 
\begin{multline}
|\nabla( f_j(y))|=\left|(\nabla(y_j-y_{j0}))\eta\left(\frac{y-y_0}{r}\right)+(y_j-y_{j0})\nabla\biggl(\eta\left(\frac{y-y_0}{r}\right)\biggr)\right|\\
=
\eta\left(\frac{y-y_0}{r}\right)+\frac{y_j-y_{j0}}{r}(\nabla\eta)\left(\frac{y-y_0}{r}\right)\leq 1+(\nabla\eta)\left(\frac{y-y_0}{r}\right)\leq C.
\nonumber
\end{multline}

Substituting in the inequality
$$
\|\varphi^{\ast}f \mid L^1_p(\Omega)\|\leq K_{p,p}(\varphi;\Omega) \|f\mid L^1_p(\widetilde{\Omega})\|
$$
the test functions 
$$
f_j(y)=(y_j-y_{j0})\eta\left(\frac{y-y_0}{r}\right),\,\,\, j=1,...,n,
$$
we have that
\begin{equation}\label{EqFunPP}
\biggl(\int\limits_{\varphi^{-1}(B(y_0,r))}|D\varphi(x)|^p\,dx\biggr)^{1/p}\leq CK_{p,p}(\varphi;\Omega)(r^n)^{1/p},
\end{equation}
where $C$ is a constant which depends on $n$ and $p$ only. Hence the homeomorphism $\varphi$ belongs to the Sobolev space $L_{p,{\loc}}^1(\Omega)$. 

Now we prove that $\varphi:\Omega\to\widetilde{\Omega}$ is a homeomorphism of finite distortion. Let $Z=\{x\in \Omega : J(x,\varphi)=0\}$. We prove that
$$
\int\limits_{Z}|D\varphi(x)|^p\,dx=0.
$$
For this aim we rewrite the integral as sum of two integrals:
$$
\int\limits_{Z}|D\varphi(x)|^p\,dx
=\int\limits_{Z\setminus S}|D\varphi(x)|^p\,dx+
\int\limits_{Z \cap S}|D\varphi(x)|^p\,dx,
$$
where $S$ is the set from the change of variables formula (\ref{chvf}) on which the homeomorphism $\varphi$ has no the Luzin $N$-property.

Because $|S|=0$, then 
$$
\int\limits_{Z\cap S}|D\varphi(x)|^p\,dx=0.
$$

Now we show that
$$
\int\limits_{Z\setminus S}|D\varphi(x)|^p\,dx=0.
$$
By the change of variable formula we have $|\varphi(Z\setminus S)|=0$. Fix $\varepsilon>0$. Then there exists a family of balls $\{B(y_i,r_i)\}$ generating a covering of the set $\varphi(Z\setminus S)$ such that the multiplicity of the covering $B(y_i,2r_i)$ is finite and $\sum\limits_i|B_i|<\varepsilon$.
Then by inequality (\ref{EqFunPP}) we obtain
$$
\int\limits_{Z\setminus S}|D\varphi(x)|^p\,dx\leq
\sum\limits_{i=1}^{\infty}\int\limits_{\varphi^{-1}(B(y_i,r_i))}|D\varphi(x)|^p\,dx
\leq C^p K^p_{p,p}(\varphi;\Omega) \sum\limits_{i=1}^{\infty}|B_i|.
$$
Since $\varepsilon$ is an arbitrary number then $\int\limits_{Z\setminus S}|D\varphi|^p\,dx=0$. Hence, $D\varphi=0$ a.~e. on $Z\setminus S$ and as a consequence $D\varphi=0$ a.e. on $Z$ and the homeomorphism $\varphi$ has finite distortion.

Now we apply to the left side of the inequality~(\ref{EqFunPP}) the change of variable formula, $B:=B(y_0,r)$:
\begin{multline}
\biggl(\int\limits_{\varphi^{-1}(B)}|D\varphi(x)|^p\,dx\biggr)^{\frac{1}{p}}
=
\biggl(\int\limits_{\varphi^{-1}(B)\setminus S}|D\varphi(x)|^p\,dx\biggr)^{\frac{1}{p}}\\
=\biggl(\int\limits_{\varphi^{-1}(B)\setminus (S\cup Z)}
|D\varphi(x)|^p\,dx\biggr)^{\frac{1}{p}}
=
\biggl(\int\limits_{\varphi^{-1}(B)\setminus (S\cup Z)}
\frac{|D\varphi(x)|^p}{|J(x,\varphi)|}|J(x,\varphi)|\,dx\biggr)^{\frac{1}{p}}\\
=\biggl(\int\limits_{B(y_0,r)\setminus \varphi(S)}
\frac{|D\varphi(\varphi^{-1}(y))|^p}{|J(\varphi^{-1}(y),\varphi)|}\,dy\biggr)^{\frac{1}{p}}
\leq CK_{p,p}(\varphi;\Omega)(r^n)^{\frac{1}{p}}.
\nonumber
\end{multline}
Hence, we have the inequality
$$
\biggl(\frac{1}{r^n}\int\limits_{B(y_0,r)\setminus \varphi(S)}
\frac{|D\varphi(\varphi^{-1}(y))|^p}{|J(\varphi^{-1}(y),\varphi)|}\,dy\biggr)^{\frac{1}{p}}
\leq CK_{p,p}(\varphi;\Omega).
$$

By using the Lebesgue theorem about differentiability of the integral by measure (see, for example \cite{RR55,VU04}) we obtain
$$
\biggl(
\frac{|D\varphi(\varphi^{-1}(y))|^p}{|J(\varphi^{-1}(y),\varphi)|}\biggr)^{\frac{1}{p}}
\leq CK_{p,p}(\varphi;\Omega)\quad\text{ for almost all}\,\,\, y\in\widetilde{\Omega}\setminus\varphi(S).
$$
Because on the set $\Omega\setminus S$ the homeomorphism $\varphi$ has the Luzin $N$-property, finally we have:
$$
\ess\sup\limits_{x\in\Omega}\left(\frac{|D\varphi(x)|^p}{|J(x,\varphi)|}\right)^{\frac{1}{p}}\leq CK_{p,p}(\varphi;\Omega)<\infty.
$$

{\it Sufficiency.}
Let $f\in L_p^1(\widetilde{\Omega})\cap C^\infty(\widetilde{\Omega})$, then the composition $f\circ\varphi$
belongs to the class $W^1_{1,\loc}(\Omega)$ and the chain rule holds \cite{Zi}. So, we have
\begin{multline}
\|\varphi^* f\vert L_p^1(\Omega)\|=\biggl(\int\limits_\Omega|\nabla (f\circ\varphi(x))|^p\,dx\biggr)^{\frac{1}{p}}
\leq
\biggl(\int\limits_\Omega(|\nabla f|(\varphi(x))|D\varphi(x)|)^p\,dx\biggr)^{\frac{1}{p}}\\
=\biggl(\int\limits_\Omega|\nabla f|^p(\varphi(x))|J(x,\varphi)|
\frac{|D\varphi(x)|)^p}{|J(x,\varphi)|}\,dx\biggr)^{\frac{1}{p}}.
\nonumber
\end{multline}
Hence
\begin{multline}
\|\varphi^*f\vert L_p^1(\Omega)\|=\biggl(\int\limits_\Omega|\nabla (f\circ\varphi(x))|^p\,dx\biggr)^{\frac{1}{p}}\\
\leq
\ess\sup\limits_{x\in\Omega}\left(\frac{|D\varphi(x)|^p}{|J(x,\varphi)|}\right)^{\frac{1}{p}}\biggl(\int\limits_\Omega|\nabla f|^p(\varphi(x))|J(x,\varphi)|\,dx\biggr)^{\frac{1}{p}}.
\nonumber
\end{multline}

The application of the change of variable formula gives the required inequality
$$
\|\varphi^{\ast}f \mid L^1_p(\Omega)\|\leq K_{p,p}(\varphi;\Omega) \|f \mid L^1_p(\widetilde{\Omega})\|
$$
for every smooth function $f\in L^1_p(\widetilde{\Omega})$.

To extend the estimate onto all functions $f\in L^1_p(\widetilde{\Omega})$, $1<q< p<\infty$, consider a sequence of smooth functions $f_k\in L^1_p(\widetilde{\Omega})$, $k=1,2,...$, such that $f_k\to f$ in $L^1_p(\widetilde{\Omega})$ and $f_k\to f$ $p$-quasi-everywhere in $\widetilde{\Omega}$ as $k\to\infty$. Since by Corollary~\ref{cor:Capacity} the preimage $\varphi^{-1}(S)$ of the set $S\subset \widetilde{\Omega}$ of $p$-capacity zero has the $p$-capacity zero, we have $\varphi^{\ast}(f_k)\to \varphi^{\ast}(f)$
$p$-quasi-everywhere in $\Omega$ as $k\to\infty$. This observation leads us to the following conclusion: Extension by continuity of the operator $\varphi^{\ast}$ $L^1_p(\widetilde{\Omega})\cap C^{\infty}(\widetilde{\Omega})$ to $L^1_p(\widetilde{\Omega})$ coincides with the composition operator $\varphi^{\ast}$, $\varphi^{\ast}(f) = f\circ\varphi$.

If $p=1$ then we should replace the capacitary characteristic of convergence with a coarser
one: if a sequence $f_n\in L^1_1(\widetilde{\Omega})$ converges to $f\in L^1_1(\widetilde{\Omega})$ in $L^1_1(\widetilde{\Omega})$ then some its subsequence converges almost everywhere. To complete the proof, it suffices to use the following property: the inverse image of a set of measure zero under the mapping $\varphi:\Omega\to\widetilde{\Omega}$ inducing the bounded operator
$\varphi^{\ast}:L^1_1(\widetilde{\Omega})\to L^1_1(\Omega)$ is a set of measure zero.

\end{proof}

The next theorem gives the characterization of composition operators from Sobolev spaces $L^1_p(\widetilde{\Omega})$ to $L^1_q(\Omega)$ for $q<p$.


\begin{thm}
\label{thm:AnalyticPQ}
Let $\Omega$ and $\widetilde{\Omega}$ be domains in $\mathbb R^n$. Then a homeomorphism $\varphi :\Omega\to \widetilde{\Omega}$ generates a bounded composition operator
$$
\varphi^{\ast} : L^1_p(\widetilde{\Omega})\to L^1_q(\Omega), \,\,\,1\leq q< p<\infty,
$$
if and only if $\varphi\in L^1_{q,\loc}(\Omega)$ has finite distortion and
$$
\biggl(\int\limits_\Omega\bigg(\frac{|D\varphi(x)|^p}{|J(x,\varphi)|}\biggr)^{\frac{q}{p-q}}~dx\biggr)^{\frac{p-q}{pq}}=K_{p,q}(\varphi;\Omega)<\infty.
$$
The norm of the operator $\|\varphi^{\ast}\|\leq K_{p,q}(\varphi;\Omega)$.
\end{thm}

\begin{proof}
{\it Necessity.} By Theorem \ref{thm:SobolevAddFun} the inequality
\begin{equation}
\label{eqphi} 
\|\varphi^{\ast} f|{L}_{q}^{1}(\Omega)\|\leq
\Phi(\widetilde{A})^{\frac{p-q}{pq}}
\|f|{L}_{p}^{1}(\widetilde{\Omega})\|,\,\,\,1\leq q<p<\infty,
\end{equation}
holds for any function $f\in {L}_{p}^{1}(\widetilde{\Omega})\cap C_0(\widetilde{A})$.

Fix a cut function
$\eta\in C_{0}^{\infty}(\mathbb R^n)$,
which is equal to one on the ball $B(0,1)$
and is equal to zero outside of the ball $B(0,2)$.
Now we consider the test functions
$$
f_j(y)=(y_j-y_{j0})\eta\left(\frac{y-y_0}{r}\right),\,\,\, j=1,...,n,
$$
where $y_j$ is the $j$-th coordinate function. We see that 
\begin{multline}
|\nabla( f_j(y))|=\left|(\nabla(y_j-y_{j0}))\eta\left(\frac{y-y_0}{r}\right)+(y_j-y_{j0})\nabla\biggl(\eta\left(\frac{y-y_0}{r}\right)\biggr)\right|\\
=
\eta\left(\frac{y-y_0}{r}\right)+\frac{y_j-y_{j0}}{r}(\nabla\eta)\left(\frac{y-y_0}{r}\right)\leq 1+(\nabla\eta)\left(\frac{y-y_0}{r}\right)\leq C.
\nonumber
\end{multline}

Substituting in the inequality (\ref{eqphi})
the test functions
$$
f_j(y)=(y_j-y_{j0})\eta\left(\frac{y-y_0}{r}\right),\,\,\, j=1,...,n,
$$
we see that
\begin{equation}\label{EqFun}
\biggl(\int\limits_{\varphi^{-1}(B(y_0,r))}|D\varphi(x)|^q\,dx\biggr)^{1/q}\leq
C\Phi(B(y_0,2r))^{\frac{p-q}{pq}}(r^n)^{1/p}.
\end{equation}
Hence the homeomorphism $\varphi$ belongs to the Sobolev space $L_{q,{\loc}}^1(\Omega)$. 

Now we prove that $\varphi$ is a homeomorphism of finite distortion.
Let $Z=\{x\in \Omega : J(x,\varphi)=0\}$. We prove that
$$
\int\limits_{Z}|D\varphi(x)|^q\,dx=0.
$$
For this aim we rewrite the integral as sum of two integrals:
$$
\int\limits_{Z}|D\varphi(x)|^q\,dx
=\int\limits_{Z\setminus S}|D\varphi(x)|^q\,dx+
\int\limits_{Z \cap S}|D\varphi(x)|^q\,dx,
$$
where $S$ is the set from the change of variables formula (\ref{chvf}) on which the homeomorphism $\varphi$ has no the Luzin $N$-property, $|S|=0$.

Hence,
$$
\int\limits_{Z\cap S}|D\varphi(x)|^q\,dx=0.
$$
Now we prove that
$$
\int\limits_{Z\setminus S}|D\varphi(x)|^q\,dx=0.
$$
By the change of variable formula (\ref{chvf}) we have that $|\varphi(Z\setminus S)|=0$. Fix $\varepsilon>0$. Then there exists a family of balls $\{B(y_i,r_i)\}$ generates a covering of the set $\varphi(Z\setminus S)$ such that the multiplicity of the covering $B(y_i,2r_i)$ is finite and $\sum\limits_i|B_i|<\varepsilon$.
Then by inequality (\ref{EqFun}) we obtain
\begin{multline}
\int\limits_{Z\setminus S}|D\varphi(x)|^q\,dx\leq
\sum\limits_{i=1}^{\infty}\int\limits_{\varphi^{-1}(B(y_i,r_i))}|D\varphi(x)|^q\,dx\\
\leq C\sum\limits_{i=1}^{\infty}\Phi(B(y_i,2r_i))^{\frac{p-q}{p}}(r_i^n)^{q/p}\\
\leq C\sum\limits_{i=1}^{\infty}\Phi(B(y_i,2r_i))^{\frac{p-q}{p}}
(\sum\limits_{i=1}^{\infty}r_i^n)^{q/p}.
\nonumber
\end{multline}
Since $\varepsilon$ is an arbitrary number then $\int\limits_{Z\setminus S}|D\varphi(x)|^q\,dx=0$. Form here we have that $D\varphi=0$ a.~e. on $Z\setminus S$ and the homeomorphism $\varphi$ has finite distortion.

Mow rewrite inequality~(\ref{EqFun}) to the form
$$
\biggl(
\int\limits_{\varphi^{-1}(B(y_0,r))}|D\varphi(x)|^q\,dx
\biggr)^{\frac{p}{p-q}}
\leq C\frac{\Phi(B(y_0,2r))}{|B(y_0,2r)|}(r^n)^{\frac{p}{p-q}}
$$
and apply to the left side of this inequality the change of variable formula, $B:=B(y_o,r)$:
\begin{multline}
\biggl(\int\limits_{\varphi^{-1}(B)}
|D\varphi(x)|^q\,dx\biggr)^{\frac{p}{p-q}}=
\biggl(\int\limits_{\varphi^{-1}(B)\setminus S}
|D\varphi(x)|^q\,dx\biggr)^{\frac{p}{p-q}}\\
=\biggl(\int\limits_{\varphi^{-1}(B)\setminus (S\cup Z)}
|D\varphi(x)|^q\,dx\biggr)^{\frac{p}{p-q}}=
\biggl(\int\limits_{\varphi^{-1}(B)\setminus (S\cup Z)}
\frac{|D\varphi(x)|^q}{|J(x,\varphi)|}|J(x,\varphi)|\,dx\biggr)^{\frac{p}{p-q}}\\
=\biggl(\int\limits_{B(y_0,r)\setminus \varphi(S)}
\frac{|D\varphi(\varphi^{-1}(y))|^q}{|J(\varphi^{-1}(y),\varphi)|}\,dy\biggr)^{\frac{p}{p-q}}
\leq C\frac{\Phi(B(y_0,2r))}{|B(y_0,2r)|}(r^n)^{\frac{p}{p-q}}.
\nonumber
\end{multline}
Hence, we obtain the inequality
$$
\biggl(\frac{1}{r^n}\int\limits_{B(y_0,r)\setminus \varphi(S)}
\frac{|D\varphi(\varphi^{-1}(y))|^q}{|J(\varphi^{-1}(y),\varphi)|}\,dy\biggr)^{\frac{p}{p-q}}
\leq C\frac{\Phi(B(y_0,2r))}{|B(y_0,2r)|}.
$$

Using the Lebesgue theorem about differentiability of the integral and properties of the volume derivative of the countable-additive set functions \cite{RR55,VU04} we obtain
$$
\biggl(
\frac{|D\varphi(\varphi^{-1}(y))|^q}{|J(\varphi^{-1}(y),\varphi)|}\biggr)^{\frac{p}{p-q}}
\leq C\Phi'(y)\quad\text{ for almost all}\,\,\, y\in\widetilde{\Omega}\setminus\varphi(S).
$$
Integrating of the last inequality on an arbitrary open bounded subset $\widetilde{U}\subset\widetilde{\Omega}$ we obtain
\begin{multline*}
\int\limits_{\widetilde{U}\setminus\varphi(S)}\biggl(
\frac{|D\varphi(\varphi^{-1}(y))|^q}{|J(\varphi^{-1}(y),\varphi)|}\biggr)^{\frac{p}{p-q}}~dy
\leq C\int\limits_{\widetilde{U}\setminus\varphi(S)}\Phi'(y)~dy\\
\leq C\int\limits_{\widetilde{U}}\Phi'(y)~dy
\leq C \Phi(\widetilde{U})\leq C  K_{p,q}^{\frac{p-q}{pq}}(\varphi;\Omega).
\end{multline*}
Since the choice of $\widetilde{U}\subset\widetilde{\Omega}$ is arbitrary, we have
$$
\int\limits_{\widetilde{\Omega}\setminus\varphi(S)}\biggl(
\frac{|D\varphi(\varphi^{-1}(y))|^q}{|J(\varphi^{-1}(y),\varphi)|}\biggr)^{\frac{p}{p-q}}~dy
\leq C K_{p,q}^{\frac{p-q}{pq}}(\varphi;\Omega).
$$
Hence
\begin{multline}
\int\limits_\Omega\left(\frac{|D\varphi(x)|^p}
{|J(x,\varphi)|}\right)^{\frac{q}{p-q}}\,dx=
\int\limits_\Omega\frac{|D\varphi(x)|^{\frac{pq}{p-q}}}
{|J(x,\varphi)|^{\frac{p}{p-q}}}|J(x,\varphi)|\,dx\\
=
\int\limits_{\widetilde{\Omega}\setminus\varphi(S)}\biggl(
\frac{|D\varphi(\varphi^{-1}(y))|^q}{|J(\varphi^{-1}(y),\varphi)|}\biggr)^{\frac{p}{p-q}}~dy
\leq C K_{p,q}^{\frac{p-q}{pq}}(\varphi;\Omega).
\nonumber
\end{multline}

{\it Sufficiency.}
Let $f\in L_p^1(\widetilde{\Omega})\cap C^\infty(\widetilde{\Omega})$,  then the composition $f\circ\varphi$
belongs to the class $W^1_{1,\loc}(\Omega)$ and the chain rule holds \cite{Zi}. So, we have
\begin{multline}
\|\varphi^* f\vert L_q^1(\Omega)\|=\biggl(\int\limits_\Omega|\nabla (f\circ\varphi)|^q\,dx\biggr)^{\frac{1}{q}}\\
\leq
\biggl(\int\limits_\Omega(|\nabla f||D\varphi(x)|)^q\,dx\biggr)^{\frac{1}{q}}
=\biggl(\int\limits_\Omega|\nabla f|^q|J(x,\varphi)|^{\frac{q}{p}}
\frac{|D\varphi(x)|)^q}{|J(x,\varphi)|^{\frac{q}{p}}}\,dx\biggr)^{\frac{1}{q}}.
\nonumber
\end{multline}
Using the H\"older inequality we obtain
\begin{multline}
\|\varphi^*f\vert L_q^1(\Omega)\|=\biggl(\int\limits_\Omega|\nabla (f\circ\varphi)|^q\,dx\biggr)^{\frac{1}{q}}\\
\le
\biggl(\int\limits_\Omega|\nabla f|^p(\varphi(x))\cdot|J(x,\varphi)|\,dx\biggr)^{\frac{1}{p}}
\cdot\biggl(\int\limits_\Omega\left(\frac{|D\varphi(x)|^p}
{|J(x,\varphi)|}\right)^{\frac{q}{p-q}}\,dx\biggr)^{\frac{p-q}{pq}}.
\nonumber
\end{multline}

Now the application of the change of variable formula gives the required inequality
$$
\|\varphi^{\ast}f \mid L^1_q(\Omega)\|\leq K \|f \mid L^1_p(\widetilde{\Omega})\|
$$
for every smooth function $f\in L^1_p(\widetilde{\Omega})$.

To extend the estimate onto all functions $f\in L^1_p(\widetilde{\Omega})$, $1<q< p<\infty$, consider a sequence of smooth functions $f_k\in L^1_p(\widetilde{\Omega})$, $k=1,2,...$, such that $f_k\to f$ in $L^1_p(\widetilde{\Omega})$ and $f_k\to f$ $p$-quasi-everywhere in $\widetilde{\Omega}$ as $k\to\infty$. Since by Corollary~\ref{cor:Capacity} the preimage $\varphi^{-1}(S)$ of the set $S\subset \widetilde{\Omega}$ of $p$-capacity zero has the $q$-capacity zero, we have $\varphi^{\ast}(f_k)\to \varphi^{\ast}(f)$
$q$-quasi-everywhere in $\Omega$ as $k\to\infty$. This observation leads us to the following conclusion: Extension by continuity of the operator $\varphi^{\ast}$ $L^1_p(\widetilde{\Omega})\cap C^{\infty}(\widetilde{\Omega})$ to $L^1_p(\widetilde{\Omega})$ coincides with the composition operator  $\varphi^{\ast}$, $\varphi^{\ast}(f) = f\circ\varphi$.

If $1 = q<p$, consider a sequence of smooth functions $f_k\in L^1_p(\widetilde{\Omega})$, $k=1,2,...$, such that $f_k\to f$ in $L^1_p(\widetilde{\Omega})$ and $f_k\to f$ $p$-quasi-everywhere in $\widetilde{\Omega}$ as $k\to\infty$. 
Then there exists a subsequence $\varphi^{\ast}(f_{n_k})$ which converges almost everywhere in $\Omega$. Hence extension by continuity of the operator $\varphi^{\ast}$ $L^1_p(\widetilde{\Omega})\cap C^{\infty}(\widetilde{\Omega})$ to $L^1_p(\widetilde{\Omega})$ coincides with the operator of substitution $\varphi^{\ast}$, $\varphi^{\ast}(f) = f\circ\varphi$.
\end{proof}

\vskip 0.3cm
In the case  $1\leq q\leq p=\infty$ we need the following lemma:

\begin{lem}
\label{lem:WeakConv}
Let a function $f: \Omega\to \mathbb R^n$ belong to the Sobolev space $L^1_{\infty}(\Omega)$. Then there exists a sequence of smooth functions $\{f_k\}$ such that the sequence $\nabla f_k$ weakly converges to $\nabla f$ in $L^{\infty}(\Omega)$.
\end{lem}

\begin{proof} Let $\Omega_{\delta} = \{x\in \Omega :\dist(x,\partial \Omega)>\delta\}$ for $\delta>0$. For $f\in L^1_{\infty}(\Omega)$, $x\in \Omega_{\delta}$, and $0<r<\delta$ a mollifier function of $\nabla f$:
$$
A_{r}\nabla f(x)=\frac{1}{{r}^n}\int\limits_{\mathbb R^n} \omega\biggl(\frac{x-z}{r}\biggr)\nabla f(z)~dz
$$
is well defined.
Then for every function $g\in C^{\infty}_{0}(\Omega)$
\begin{multline}
\biggl|\int\limits_{\Omega_{\delta}} A_r \nabla f(x)g(x)dx - \int\limits_{\Omega_{\delta}} \nabla f(x)g(x)dx\biggr|
=\biggl|\int\limits_{\Omega_{\delta}} (A_r \nabla f(x) - \nabla f(x))g(x)~dx\biggr|
\\
=\biggl|\int\limits_{\Omega_{\delta}}\biggl( \frac{1}{r^n}\int\limits_\Omega\omega\biggl(\frac{x-y}{r}\biggr)\nabla f(x)dy - \nabla f(x)\biggr)g(x)~dx\biggr|\\
=\biggl|\int\limits_{\Omega_{\delta}}\biggl( \int\limits_{B(0,1)}\nabla f(x-r z)\omega(z)dz - \nabla f(x)\biggr)g(x)~dx\biggr|\\
=\biggl|\int\limits_{\Omega_{\delta}}\biggl( \int\limits_{B(0,1)}(f(x-r z)  - f(x))\omega(z)dz \biggr)g(x)dx\biggr|.
\nonumber
\end{multline}


Thus, we have:
\begin{multline}
\biggl|\int\limits_{\Omega_{\delta}} g(x)A_r \nabla f(x)~dx - \int\limits_{\Omega_{\delta}}g(x)\nabla f(x)~dx\biggr|\\
\leq
\int\limits_{B(0,1)}\biggr( \int\limits_{\Omega_{\delta}}|\nabla f(x-r z) - \nabla f(x)||g(x)|~dx\biggl)\omega(z)~dz.
\nonumber
\end{multline}

Now, if $\nabla f$ is a continuous vector-function, then 
$$
\biggl|\int\limits_{\Omega_{\delta}}g(x) A_r \nabla f(x)~dx - \int\limits_{\Omega_{\delta}} g(x) \nabla f(x)~dx\biggr|
\to 0\,\,\,\text{as}\,\,\,r\to 0.
$$

Since $f\in L^1_{\infty}(\Omega)$ then for every $\varepsilon>0$ by the Luzin theorem $\nabla f=f_1+f_2$ where $f_1$ is a continuous vector-function and $f_2$ such that
$$
\int\limits_{\Omega_{\delta}}|f_2(x-r z) - f_2(x)||g(x)|~dx<\frac{\varepsilon}{2} ~.
$$


Choose by continuity of $f_1$ a number $r>0$ such that
$$
\int\limits_{\Omega_{\delta}}|f_1(x- r z) - f_1(x)||g(x)|~dx<\frac{\varepsilon}{2}~.
$$
Then
\begin{equation}
\biggl|\int\limits_{\Omega_{\delta}} g(x) A_r \nabla f(x)~dx - \int\limits_{\Omega_{\delta}} g(x) \nabla f(x)~dx\biggr|<\frac{\varepsilon}{2}+\frac{\varepsilon}{2}<\varepsilon~.
\end{equation}


So, since the space $C^{\infty}_{0}(\Omega)$ is dense in the  conjugate space to $L^{\infty}(\Omega)$ and $\|\nabla (A_r f)\|$ converges to $\|\nabla f\|$ (see \cite{Bu98})
then 
$\nabla (A_r f)$ weakly converges to $\nabla f$ in $L^{\infty}(\Omega_{\delta})$ (see, for example, Ch.~5 from \cite{Yo}).

The weak convergence of gradients in $L^{\infty}(\Omega)$ is proved in the standard manner.
Indeed, choose a sequence of open sets $\Omega_j \Subset \Omega_{j+1}\Subset \Omega$, $j\geq 1$ compactly embedded one into another such that $\cup_{j}\Omega_j=\Omega$. Let $\Psi$ be a partition of unity on $\Omega$, subordinate to the covering $\Omega_{j+1}\setminus\overline{\Omega}_{j-1}$, see e.~g. Lemma~2.3.1 in \cite{Zi}. Let $\psi_j$ denote the (finite) sum of those $\psi\in \Psi$ for which 
$\supp\psi\subset \Omega_{j+1}\setminus\overline{\Omega}_{j-1}$. Thus $\psi_j\in C_0^{\infty}(\Omega_{j+1}\setminus\overline{\Omega}_{j-1})$ and $\sum_j \psi_j\equiv 1$ in $\Omega$.

Given $g\in C^{\infty}_{0}(\Omega)$, $\gamma>0$ and $j=1,2,\dots$. Then  by (1) for $\varepsilon=\gamma 2^{-j}$  
$$
\biggl|\int\limits_{\Omega} g(x) \nabla \eta_j(x)~dx - \int\limits_{\Omega} g(x) \nabla\bigl(\psi_j(x)f(x)\bigr)~dx\biggr|\leq \gamma 2^{-j}
$$
where $\eta_j=A_{r_j}\psi_j f$ for small enough $r_j$.

Then, setting $\eta=\sum_j \eta_j\in C^{\infty}(\Omega)$ we have that
\begin{multline}
\biggl|\int\limits_{\Omega} g(x) \nabla \eta(x)~dx - \int\limits_{\Omega} g(x) \nabla f(x)~dx\biggr|\\
=
\biggl|\int\limits_{\Omega} g(x)\sum_j \nabla \eta_j(x)~dx - \int\limits_{\Omega} g(x) \nabla\bigl(\sum_j\psi_j(x)f(x)\bigr)~dx\biggr|\\
\leq\sum_j\biggl|\int\limits_{\Omega} g(x) \nabla \eta_j(x)~dx - \int\limits_{\Omega} g(x) \nabla\bigl(\psi_j(x)f(x)\bigr)~dx\biggr|<\gamma.
\nonumber
\end{multline}

Since the space $L^1(\Omega)$ is separable, it is possible to find a sequence $g_i$, say in $C_0^{\infty}(\Omega)$, which is dense in the conjugate space to $L^{\infty}(\Omega)$, see e.~g. \cite{Yo}. Choosing for $\gamma=\gamma_k=2^{-k}$, $r_j=r_j(k)$ in the above construction the same for $g=g_i$, $i=1,2,...,k$, we obtain that $\nabla f_k\to \nabla f$ weakly in $L^{\infty}(\Omega)$ where 
$$
f_k=\sum\limits_{j=1}^{\infty} A_{r_j(k)}\psi_jf,\,\,\, k=1,2,...~.
$$
\end{proof}

\begin{thm}
\label{thm:AnalyticInfP}
Let $\varphi:\Omega\to \widetilde{\Omega}$ be a homeomorphism between two domains $\Omega,\widetilde{\Omega} \subset \mathbb R^n$. 
Then $\varphi$ belongs to the Sobolev space $L^1_p(\Omega)$, $1\leq p\leq \infty$ if and only if the composition operator
$$
\varphi^{\ast}: L^1_{\infty}(\widetilde{\Omega})\to L^1_{p}(\Omega)
$$
is bounded.
\end{thm}

\begin{proof}
{\it Necessity.} 
Since the composition operator $\varphi^{\ast}$ is bounded, substituting in the inequality
$$
\|\varphi^{\ast}f \mid L^1_p(\Omega)\|\leq \|\varphi^{\ast}\|\|f\mid L^1_{\infty}(\widetilde{\Omega})\|
$$
the coordinate functions $f_j=y_j$, $j=1,...,n$, we obtain that
$$
\|\varphi_j\mid L^1_p(\Omega)\|\leq \|\varphi^{\ast}\|, \,\,\,j=1,...,n,
$$
because $\|y_j\mid L^1_{\infty}(\widetilde{\Omega})\|=1$. Hence the mapping $\varphi$ belongs to the Sobolev space $L^1_p(\Omega)$.

{\it Sufficiency.} Let $f$ be a smooth function of the class $L^1_{\infty}(\widetilde{\Omega})$. Then the composition $f\circ\varphi$ belongs to the Sobolev space $W^1_{1,\loc}(\Omega)$ and
\begin{multline}
\|\varphi^{\ast}f\mid L^1_p(\Omega)\|=\biggl(\int\limits_\Omega|\nabla(f\circ\varphi)|^p~dx\biggr)^{\frac{1}{p}}\leq
\biggl(\int\limits_\Omega|D\varphi(x)|^p|\nabla f|^p(\varphi(x))~dx\biggr)^{\frac{1}{p}}\\
\leq \biggl(\int\limits_\Omega|D\varphi(x)|^p~dx\biggr)^{\frac{1}{p}} \|f\mid L^1_{\infty}(\widetilde{\Omega})\|.
\nonumber
\end{multline}

For an arbitrary function $f\in L^1_{\infty}(\widetilde{\Omega})$ we find by Lemma~\ref{lem:WeakConv} a sequence of smooth functions $f_k\in L^1_{\infty}(\widetilde{\Omega})$ such that $\nabla f_k\to \nabla f$ weakly in $L^{\infty}(\widetilde{\Omega})$ and 
$$
\|f_k\mid L^1_{\infty}(\widetilde{\Omega})\|\to\|f\mid L^1_{\infty}(\widetilde{\Omega})\|,
$$
see e.~g. \cite{Bu98}.
Note that 
$$
\|f\mid L^1_{\infty}(\widetilde{\Omega})\|\leq\liminf\limits_{k\to\infty} \|f_k\mid L^1_{\infty}(\widetilde{\Omega})\|,
$$
see e.~g. Lemma II. 3.27 in \cite{DSh58}.
Then $\nabla(\varphi^{\ast}f_k)\to \nabla(\varphi^{\ast}f)$ weakly in $L^p(\Omega)$ and
$$
\|\nabla(\varphi^{\ast}f) \mid L^p(\Omega)\|\leq \liminf\limits_{k\to \infty}\|\nabla(\varphi^{\ast}f_k)\mid L^{p}(\Omega)\|=\liminf\limits_{k\to \infty}\|\varphi^{\ast}f_k\mid L^1_p(\Omega)\|.
$$
Hence, tending to the limit in the inequality
$$
\|\varphi^{\ast}f_k\mid L^1_p(\Omega)\|\leq K \|f_k\mid L^1_{\infty}(\widetilde{\Omega})\|
$$
we obtain that
$$
\|\varphi^{\ast}f\mid L^1_p(\Omega)\|\leq K \|f\mid L^1_{\infty}(\widetilde{\Omega})\|.
$$
\end{proof}


\subsection{Mappings of bounded $(p,q)$-distortion}

Let us define the "normalized" $p$-dilatation of a Sobolev homeomorphism $\varphi:\Omega\to\widetilde{\Omega}$ at a point $x\in\Omega$ as
$$
K_p(x)=\inf \{k(x): |D\varphi(x)|\leq k(x) |J(x,\varphi|^{\frac{1}{p}},\,\,x\in\Omega \}.
$$
In the case $p=n$ we have the "normalized" conformal dilatation and in the case $p\ne n$ the  equivalent $p$-dilatation was introduced in \cite{Ge69}. 
Then $\varphi:\Omega\to\widetilde{\Omega}$ is called a mapping of bounded weak $(p,q)$-distortion \cite{UV10} or a weak $(p,q)$-quasiconformal mapping \cite{GGR95,VU98} if $\varphi\in W^1_{q,\loc}(\Omega)$, has finite distortion,
and 
\[
K_{p,q}(\varphi;\Omega)=\|K_p \mid L_{\kappa}(\Omega)\|<\infty,
\]
where $1/q-1/p=1/{\kappa}$ ($\kappa=\infty$, if $p=q$).


Let us consider as an example mappings of bounded weak $(p,q)$-distortion of Lipschitz domains onto anisotropic H\"older singular domains (introduced in \cite{GGu94}).
Let $H_g$ be a domain with anisotropic H\"older $\gamma$-singularities:
$$
H_g=\{ x\in\mathbb R^n : 0<x_n<1, 0<x_i<g_i(x_n),
\,i=1,2,\dots,n-1\}.
$$
Here $g_i(\tau)=\tau^{\gamma_i}$, $\gamma_i\geq 1$, $0\leq\tau\leq 1$ are H\"older functions and for the function $G=\prod_{i=1}^{n-1}g_i$ denote by
$$
\gamma=\frac{\log G(\tau)}{\log \tau}+1.
$$
It is evident that $\gamma\geq n$. In the case $g_1=g_2=\dots=g_{n-1}$ we will say that domain $H_g$ is a domain with $\sigma$-H\"older singularity, $\sigma=(\gamma-1)/(n-1)$.
For $g_1(\tau)=g_2(\tau)=\dots=g_{n-1}(\tau)=\tau$ we will use notation $H_1$ instead
of $H_g$.

Define the mapping $\varphi_a: H_1\to H_g$, $a>0$, by
$$
\varphi_a(x)=\left(\frac{x_1}{x_n}g^a_1(x_n),\dots,\frac{x_{n-1}}{x_n}g^a_{n-1}(x_n),x_n^a\right).
$$

\begin{thm}
\label{lemhol}
Let $(n-p)/(\gamma-p)<a<p(n-q)/q(\gamma-p)$. Then the mapping $\varphi_a: H_1\to H_g$, 
is a weak $(p,q)$-quasiconformal mapping, $1<q<p<\gamma$, from the Lipschitz convex domain $H_1$ onto the "cusp"' domain $H_g$ with
$$
K_{p,q}(H_1)\leq a^{-\frac{1}{p}}\sqrt{\sum_{i=1}^{n-1}(a\gamma_i-1)^2+n-1+a^2}\,.
$$
\end{thm}

\begin{proof}
By simple calculations 
$$
\frac{\partial(\varphi_a)_i}{\partial
x_i}=\frac{g^a_i(x_n)}{x_n},\quad
\frac{\partial(\varphi_a)_i}{\partial
x_n}=\frac{-x_ig^a_i(x_n)}{x_n^{2}}+\frac{ax_ig^{a-1}_i(x_n)}{x_n}g'_i(x_n)
\quad\text{and}\quad\frac{\partial(\varphi_a)_n}{\partial
x_n}=ax_n^{a-1}
$$
for any $i=1,...,n-1$. Hence $J(x,\varphi_a)=ax_n^{a-n}G^a(x_n)=ax_n^{a\gamma-n}$, $J(x,\varphi_a)\leq a$ for $a>1$ and
\begin{multline}\label{maps}
D\varphi_a (x)=
\left(\begin{array}{cccc}
x_n^{a\gamma_1-1} & 0 & ... & (a\gamma_1-1)x_1x_n^{a\gamma_1-2}\\
0 & x_n^{a\gamma_2-1} & ... & (a\gamma_2-1)x_2x_n^{a\gamma_2-2}\\
... & ... & ... & ...\\
0 & 0 & ... & ax_n^{a-1}
\end{array} \right)\\
=
x_n^{a-1}\left(\begin{array}{cccc}
x_n^{a\gamma_1-a} & 0 & ... & (a\gamma_1-1)\frac{x_1}{x_n}x_n^{a(\gamma_1-1)}\\
0 & x_n^{a\gamma_2-a} & ... & (a\gamma_2-1)\frac{x_2}{x_n}x_n^{a(\gamma_2-1)}\\
... & ... & ... & ...\\
0 & 0 & ... & a
\end{array} \right).
\end{multline}

Because $0<x_n<1$ and $x_1/x_n<1$ we have that
\begin{equation}
|D\varphi_a(x)|\leq x_n^{a-1}\sqrt{\sum_{i=1}^{n-1}(a\gamma_i-1)^2+n-1+a^2}.
\nonumber
\end{equation}

Then 
\begin{multline*}
K_{p,q}(H_1)=\left(\int\limits_{H_1}\left(\frac{|D\varphi_a(x)|^p}{J(x,\varphi_a)}\right)^{\frac{q}{p-q}}~dx\right)^{\frac{p-q}{pq}}
\\
\leq  \frac{\sqrt{\sum_{i=1}^{n-1}(a\gamma_i-1)^2+n-1+a^2}}{\sqrt[p]{a}}\left(\int\limits_{H_1}x_n^{\frac{(p(a-1)-(a\gamma-n))q}{p-q}}~dx\right)^{\frac{p-q}{pq}}\\
=
\frac{\sqrt{\sum_{i=1}^{n-1}(a\gamma_i-1)^2+n-1+a^2}}{\sqrt[p]{a}} \left(\int\limits_0^1\int\limits_0^{x_n}...\int\limits_0^{x_n}x_n^{\frac{(p(a-1)-(a\gamma-n))q}{p-q}}~dx_1...dx_n\right)^{\frac{p-q}{pq}}\\=
\frac{\sqrt{\sum_{i=1}^{n-1}(a\gamma_i-1)^2+n-1+a^2}}{\sqrt[p]{a}}\left(\int\limits_0^1 x_n^{\frac{(p(a-1)-(a\gamma-n))q}{p-q}+n-1}~dx_n\right)^{\frac{p-q}{pq}}\\
\leq \frac{\sqrt{\sum_{i=1}^{n-1}(a\gamma_i-1)^2+n-1+a^2}}{\sqrt[p]{a}},
\end{multline*}
if $(p+a\gamma -pa)q<np$ or that equivalent $a<p(n-q)/q(\gamma-p)$. 

Now we check that $1<q<np/(p+a\gamma-pa)<p$. The inequality $1<np/(p+a\gamma-pa)$ implies $a<(np-p)/\gamma-p$, but
$$
\frac{np-p}{\gamma-p}<\frac{p(n-q)}{q(\gamma-p)}, \,\,\text{if}\,\,q>1.
$$
The inequality $np/(p+a\gamma-pa)<p$ implies $a>(n-p)/(\gamma-p)$. So, we have that $a\in \left((n-p)/(\gamma-p),p(n-q)/q(\gamma-p)\right)$.
\end{proof}


\section{On the weak inverse mapping theorem for Sobolev homeomorphisms}

In this section  we consider weak regularity properties of mappings inverse to homeomorphisms that generate bounded composition operators on Sobolev spaces. On this base we prove that mappings inverse to homeomorphisms generating bounded composition operators on Sobolev spaces
$$
\varphi^{\ast}: L^1_p(\widetilde{\Omega})\to L^1_q(\Omega),\,\,\,n-1\leq q\leq p\leq\infty,
$$
generate bounded composition operators also, see \cite{U93} (the limit case in \cite{GU10-2}).

Let us recall that the ring condenser $R$ (see, for example, \cite{Ge69,MRV69}) in a domain $\Omega\subset \mathbb R^n$ is the pair $R=(F;G)$ of sets $F\subset G\subset\Omega$, where $F$ is a connected compact set and $G$ is a connected open set.  
The $p$-capacity of the ring $R=(F;G)$ is defined by
$$
\cp_p(F;G) =\inf\|f\vert L_p^1(\Omega)\|^p,
$$
where the greatest lower bound is taken over all continuous functions $f\in L_p^1(\Omega)$ with a compact support in $G$ and such that $f\geq 1$ on $F$. Such functions are called admissible functions for the ring $R=(F;G)$. Recall the capacity estimate \cite{K86}:

\begin{lem} \label{lemcap} Let $E$ be a connected closed subset of
an open bounded set $G\subset\mathbb R^n, n\geq 2$, and $n-1<p<\infty$. Then
$$
\cp_p^{n-1}(E,G)\geq c\frac{(\diam E)^p}{|G|^{p-n+1}},
$$
where a constant $c$ depends on $n$ and $p$ only.
\end{lem}

By using this lemma we prove weak differentiability of mappings inverse to homeomorphisms generating bounded composition operators on Sobolev spaces.

\begin{thm}\label{reginv}
Let a homeomorphism $\varphi :\Omega\to \widetilde{\Omega}$
generates a bounded composition operator
$$
\varphi^{\ast}: L^1_p(\widetilde{\Omega})\to L^1_q(\Omega),\,\,\,n-1<q\leq p<\infty.
$$
Then the inverse mapping $\varphi^{-1} :\widetilde{\Omega}\to \Omega$
is $\ACL$-mapping, differentiable almost everywhere in
$\widetilde{\Omega}$.
\end{thm}


\begin{proof}
Consider an arbitrary $n$-dimensional cube $P,\,\overline{P}\subset \widetilde{\Omega}$, with edges parallel to coordinate axes. We show that $\varphi^{-1}$ is absolutely continuous on almost all section of $P$ by lines parallel to axis  $y_n$.
Let $P_0$ be a projection of $P$ to subspace $y_n=0$ and $I$ be a projection of $P$ to coordinate axis $y_n$. Then $P=P_0\times I$. 

The monotone countable-additive set function $\Phi$ by the rule $\Phi(A,P)=\Phi(A\times I)$
generates a monotone countable-additive set function defined on open subsets $A\subset P_0$. It is known that for almost all points $z\in P_0$ the value
$$
\overline{\Phi}'(z,P)=\varlimsup\limits_{r\to 0}
[\Phi(B^{n-1}(z,r),P)/|B^{n-1}(z,r)|]
$$
is finite, here $B^{n-1}$ is $(n-1)$-dimensional ball with the center at the point $z$ radius $r>0$ and
$|B^{n-1}(z,r)|$ is $(n-1)$-dimensional measure.

Similar, the  $n$-dimensional Lebesgue measure $\Psi(E)=|\varphi^{-1}(E)|$,
where $E$ is an open set in $\widetilde{\Omega}$, is the quasiadditive set function generates quasiadditive function $\Psi(A,P)=\Psi(A\times I)$
of open set $A\subset P_0$,
and $\overline{\Psi}'(z,P)<\infty$ for almost all $z\in P_0$.

Fix arbitrary point $z\in P_0$, in which the derivatives $\overline{\Phi}{}'(z,P)$ and $\overline{\Psi}'(z,P)$ are finite.
On the section $I_z=\{z\}\times I$ of the cube $P$ we take mutually disjoint intervals $\gamma_1,\dots ,\gamma_k$
with lengths $l(\gamma_1),\dots ,l(\gamma_k)$ correspondingly. Denoting by $R_i$ the set of points remote of $\gamma_i$ on distance less then given $r>0$,
consider condensers $(\gamma_i,F_i),\,F_i=
\widetilde{\Omega}\backslash(\gamma _i\cup R_i)$. Let $r>0$ is chosen such that the sets $R_1,\dots ,R_k$ are mutually disjoint, $R_i\subset P$ and
$\omega_nr<\omega_{n-1}l(\gamma_i),\,i=1,\dots ,k$, where $\omega_n$ and
$\omega_{n-1}$ are volume of the unit ball and area of its surface in $\mathbb R^n$ correspondingly. 

Now applying the inequality
$$
\cp_{q}^{1/q}(\varphi^{-1}(F_0),\varphi^{-1}(F_1);\Omega)
\leq\Phi(\widetilde{\Omega}\setminus(F_0\cup F_1))^{\frac{p-q}{pq}}
\cp_{p}^{1/p}(F_0,F_1;\widetilde{\Omega})
$$
(we put $\Phi(\widetilde{\Omega}\setminus(F_0\cup F_1))^{\frac{p-q}{pq}}=K_{p,p}(\varphi;\Omega)$ in the case $p=q$)
the upper capacity estimates \cite{GResh}
$$
{{\cp}}_p(\gamma_i,F_i;\widetilde{\Omega})\le\frac{|R_i|}{r^p}
\le 2\omega_{n-1}b_ir^{n-1-p},
$$
and the lower capacity estimates of Lemma~\ref{lemcap}
$$
{{\cp}}_q(\varphi^{-1}(\gamma_i),\varphi^{-1}(F_i);\Omega)\ge
c\frac{(\diam\varphi^{-1}(\gamma_i))^{q/(n-1)}}
{|\varphi^{-1}(R_i)|^{(1-n+q)/(n-1)}},
$$
we obtain
$$
\diam\varphi^{-1}(\Delta_i)
\le c_1\cdot r^{(n-1-p)(n-1)/p}
\Phi(R_i)^{(p-q)(n-1)/pq}\Psi(R_i)^{(1-n+q)/q}b_i^{(n-1)/p}.
$$
Summing by $i=1,\dots ,k$, applying the H\"older inequality and using the definition of the countable-additive set function we have
\begin{multline}
\sum\limits_{i=1}^{k}\diam\varphi^{-1}(\Delta_i)\le c_1
\left(\frac{\Phi(B^{n-1}(z,r),P)}{|B^{n-1}(z,r)|}\right)^
{((p-q)/(n-1))/pq}\\
\times\left(\frac{\Psi(B^{n-1}(z,r,P)}{|B^{n-1}(z,r)|}\right)^
{(1-n+q)/q}\left(\sum\limits_{i=1}^{k}l(\gamma_i)\right)^{(n-1)/p}.
\nonumber
\end{multline}
Passing $r\to 0$, we obtain
$$
\sum\limits_{i=1}^k\diam\varphi^{-1}(\gamma_i)\le c_2
\left(\sum\limits_{i=1}^k l(\gamma_i)\right)^{(n-1)/p},
$$
from here $\varphi^{-1}\in {\ACL}(\widetilde{\Omega})$.

Now we show that $\varphi^{-1}$ is differentiable almost everywhere in $\widetilde{\Omega}$.
For every point $y\in \widetilde{\Omega}$ we consider the ring condenser
$(F_0,F_1)$, where
$F_0=\{z\vert\,|y-z|\le r\},\ F_1=\{z\vert\,|y-z|\ge 2r\}$.
Then ${{\cp}}_p(F_0,F_1;\widetilde{\Omega})=c\cdot r^{n-p}$, where $c$ is some constant.
Using the lower capacity estimates we obtain
$$
\left(\frac{\diam\varphi^{-1}(F_0)}r\right)^{pq}
\le c\left(\frac{\Phi(B)}{|B|}\right)^{(n-1)(p-q)}
\left(\frac{\Psi(B)}{|B|}\right)^{pq+p-np}.
$$
Here, as before, $\Psi(B)=|\varphi^{-1}(B)|$,
$B=\widetilde{\Omega}\backslash  F_1$. Passing to the limit while $r\to 0$, for almost all points
$y\in \widetilde{\Omega}$ we obtain
$$
L(y,\varphi^{-1})\le c(\overline{\Phi}'(y))^{(n-1)(p-q)/pq}
(\overline{\Psi}'(y))^{(1-n+q)/q},
$$
where
$$
L(y,\varphi^{-1})=\varlimsup_{r\to 0}(\max_{|y-z|=r}|\varphi^{-1}(y)-
\varphi^{-1}(z)|/r),
$$
and values $\overline{\Phi}'=\varlimsup\limits_{r\to 0}(\Phi(B)/|B|)$
and $\overline{\Psi}'(y)=\varlimsup\limits_{r\to 0}(\Psi(B)/|B|)$
are finite almost everywhere in $\widetilde{\Omega}$. By using Stepanov's theorem \cite{Fe69} we obtain that $\varphi^{-1}$
is differentiable almost everywhere in $\widetilde{\Omega}$.
\end{proof}

On the base of this theorem we obtain the following duality property of composition operators on Sobolev spaces \cite{U93}. 

\begin{thm}
\label{CompDual} Let a homeomorphism $\varphi:\Omega\to\widetilde{\Omega}$,  $\Omega,\widetilde{\Omega}\subset\mathbb R^n$, induces a bounded composition
operator 
\[
\varphi^{\ast}:L^1_p(\widetilde{\Omega})\to L^1_q(\Omega),\,\,\,n-1<q\leq p<\infty.
\]
Then the inverse mapping $\varphi^{-1}:\widetilde{\Omega}\to\Omega$ induces a bounded composition operator 
\[
\left(\varphi^{-1}\right)^{\ast}:L^1_{q'}(\Omega)\to L^1_{p'}(\widetilde{\Omega}),\,\,\,n-1<p'\leq q'<\infty,
\]
where $p'=p/(p-n+1)$ and $q'=q/(q-n+1)$.
\end{thm}

\begin{proof}
By Theorem~\ref{reginv} we have that $\varphi^{-1}\in W^1_{1,\loc}(\widetilde{\Omega})$. Then by \cite{P93} (see, also, \cite{CHM,GU10-1,HKM06}) we can define
$$
|D\varphi^{-1}(y)|=
\begin{cases}
\biggl(\frac{|\adj D\varphi|(x)}{|J(x,\varphi)|}\biggr)_{x=\varphi^{-1}(y)}
& \text{if}\quad x\in \Omega\setminus \left(S\cup Z\right),
\\
\,\,0 & \text{otherwise}.
\end{cases}
$$

Hence 
\[
|D\varphi^{-1}(y)|\leq\frac{|D\varphi(x)|^{n-1}}{|J(x,\varphi)|},
\]
for almost all $x\in \Omega\setminus \left(S\cup Z\right)$, $y=\varphi(x)\in \widetilde{\Omega}\setminus \varphi\left(S\cup Z\right)$,
and  
$$
|D\varphi^{-1}(y)|=0\,\,\text{for almost all}\,\,y\in \varphi(S).
$$

Now, taking into account that
\[
\frac{q'p'}{q'-p'}=\frac{pq}{(p-q)(n-1)}
\]
we obtain 
\[
\int\limits _{\widetilde{\Omega}}\left(\frac{|D\varphi^{-1}(y)|^{q'}}{|J(y,\varphi^{-1})|}\right)^{p'/(q'-p')}\, dy\le\int\limits _{\Omega}\left(\frac{|D\varphi(x)|^{p}}{|J(x,\varphi)|}\right)^{q/(p-q)}\, dx
\]
(in the case $p=q$ we have $p'=q'$ and $L_{\infty}$-norms instead integrals) and by Theorem~\ref{thm:AnalyticPP} and Theorem~\ref{thm:AnalyticPQ} we have a bounded composition operator
\[
\left(\varphi^{-1}\right)^{\ast}:L^1_{q'}(\Omega)\to L^1_{p'}(\widetilde{\Omega}),\,\,\,n-1<p'\leq q'<\infty.
\]

\end{proof}

\begin{rem}
In the case $n=2$ we have $p'=p/(p-1)$, $q'=q/(q-1)$ and $p''=p$, $q''=q$. Hence the homeomorphism $\varphi:\Omega\to\widetilde{\Omega}$,  $\Omega,\widetilde{\Omega}\subset\mathbb R^n$, induces a bounded composition
operator 
\[
\varphi^{\ast}:L^1_p(\widetilde{\Omega})\to L^1_q(\Omega),\,\,\,1<q\leq p<\infty,
\]
if and only if the inverse mapping $\varphi^{-1}:\widetilde{\Omega}\to\Omega$ induces a bounded composition operator 
\[
\left(\varphi^{-1}\right)^{\ast}:L^1_{q'}(\Omega)\to L^1_{p'}(\widetilde{\Omega}),\,\,\,1<p'\leq q'<\infty.
\]

In the case $n\ne 2$ we have 
$$
p'' =\left(p'\right)'=\frac{p}{(n-1)^2-p(n-2)}\ne p, \,\,\text{if}\,\,p'>n-1,
$$
$$
q'' =\left(q'\right)'=\frac{q}{(n-1)^2-q(n-2)}\ne q, \,\,\text{if}\,\,q'>n-1,
$$
and this case is more complicated.
\end{rem}

Weak regularity of the inverse mapping in the case $q=n-1$ is more complicated and requires additional assumptions. 


\begin{thm}Let a homeomorphism $\varphi : \Omega\to \widetilde{\Omega}$
between two domains $\Omega$ and $\widetilde{\Omega} \subset \mathbb R^n $, $n\geq 2$, generate by the composition rule $\varphi^{\ast}f=f\circ\varphi$ a bounded operator 
$$
\varphi^{\ast}: L^1_{\infty}(\widetilde{\Omega})\to L^1_{n-1}(\Omega),
$$
possess the Luzin $N$-property and have finite distortion. Then the inverse mapping $\varphi^{-1}: \widetilde{\Omega}\to \Omega$ generates the composition operator
$$
(\varphi^{-1})^{\ast}: L^1_{\infty}(\Omega)\to L^1_{1}(\widetilde{\Omega})
$$
and belongs to the Sobolev space $L^1_1(\widetilde{\Omega})$.
\end{thm}

\begin{proof}
Let $f\in L^1_{\infty}(\widetilde{\Omega})$. Then the composition $f\circ\varphi$ is weakly differentiable in $\Omega$.
Since $\varphi$ is the mapping of finite distortion, $D\varphi(x)=0$ for a.~e. $x\in \Omega$ where $J(x,\varphi)=0$. Hence we may define $\adj D\varphi(x)=0$ at such points.
Then  
$$
|J(x,\varphi)||\nabla f|(\varphi(x))\leq |\nabla(f\circ\varphi)|(x)\adj D\varphi(x) \,\,\,\text{for almost all}\,\,\,x\in \Omega,
$$
because 
$$
\min\limits_{|h|=1}|D\varphi(x)\cdot h|=\frac{1}{\max\limits_{|h|=1}|(D\varphi(x))^{-1}\cdot h|}
$$
and
$$
(D\varphi(x))^{-1}=J^{-1}(x,\varphi)\adj D\varphi(x)
$$
if $J(x,\varphi)\ne 0$.
Hence
\begin{multline}
\|f\mid L^1_1(\widetilde{\Omega})\|=\int\limits_{\widetilde{\Omega}}|\nabla f|(y)~dy=\int\limits_{\Omega}|\nabla f|(\varphi(x))|J(x,\varphi)|~dx\\
\leq
\int\limits_{\Omega}|\nabla (f\circ\varphi)|(\varphi(x))|\adj D\varphi(x)|~dx
\leq 
\int\limits_{\Omega}|\nabla (f\circ\varphi)|(\varphi(x))|D\varphi|^{n-1}(x)~dx\\
\leq \|\varphi\mid L^1_{n-1}(\Omega)\|^{n-1}\cdot \|\varphi^{\ast}f\mid L^1_{\infty}(\Omega)\|
\nonumber
\end{multline}
because by Theorem~\ref{thm:AnalyticInfP} the mapping $\varphi$ belongs to $L^1_{n-1}(\Omega)$.
So, we have the lower estimate for the composition operator
$$
\|f\mid L^1_1(\widetilde{\Omega})\|\leq \|\varphi\mid L^1_{n-1}(\Omega)\|^{n-1}\cdot \|\varphi^{\ast}f\mid L^1_{\infty}(\Omega)\|
$$
if a function $\varphi^{\ast}f$ belongs to $L^1_{n-1}(\Omega)$.
Therefore, the inverse operator $(\varphi^{\ast})^{-1}=(\varphi^{-1})^{\ast}$ is the bounded operator 
$$
(\varphi^{-1})^{\ast} : L^1_{\infty}(\Omega)\cap L^1_{n-1}(\Omega)\to L^1_1(\widetilde{\Omega}).
$$
Fix a cut function $\eta\in C_0^{\infty}(B(0,2))$ such that $\eta=1$ on $B(0,1)$. Substituting in the inequality
$$
\|(\varphi^{-1})^{\ast}g\mid L^1_1(\widetilde{\Omega})\|\leq K \|g\mid L^1_{\infty}(\Omega)\|
$$
the test functions $g_j=(x_j-x_{0j})\eta\bigl(\frac{x-x_0}{r}\bigr)$, $x_0\in \Omega$, $r<\dist(x_0,\partial \Omega)$ we obtain that
$\varphi^{-1}$ belongs to the Sobolev space $L^1_{1,\loc}(\widetilde{\Omega})$.
Now, using the estimate
$$
\int\limits_{\widetilde{\Omega}}|D\varphi^{-1}(y)|~dy\leq\int\limits_{\widetilde{\Omega}}\frac{|D\varphi(\varphi^{-1}(y))|^{n-1}}{|J(\varphi^{-1}(y),\varphi)|}~dy=
\int\limits_{\Omega}|D\varphi(x)|^{n-1}~dx<+\infty
$$
we conclude that $\varphi^{-1}$ belongs to $L^1_1(\widetilde{\Omega})$ and by Theorem~\ref{thm:AnalyticInfP} generates a bounded composition operator from $L^1_{\infty}(\Omega)$ to $L^1_1(\widetilde{\Omega})$.
\end{proof}



\section{Metric properties of weak $(p,q)$-quasiconformal mappings}

To conclude this review we consider metric properties of weak $(p,q)$-quasiconformal mappings.

\begin{thm}
\label{Hol}
Let $\varphi:\Omega\to\widetilde{\Omega}$ be a weak $(p,q)$-quasiconformal mapping, $n<q\leq p<\infty$. Then for any points $y_1,y_2\in\widetilde{\Omega}$, such that $|y_1-y_2|=r<\dist(y_1,\partial\widetilde{\Omega})$ the following inequality
$$
|y_1-y_2|^{\frac{p-n}{p}}\leq c(n,p,q)\Phi\left((B(y_1,r)\right)^{\frac{p-q}{pq}}|\varphi^{-1}(y_1)-\varphi^{-1}(y_2)|^{\frac{q-n}{q}}
$$
holds.
\end{thm}

\begin{proof}
Let points $y_1,y_2\in\widetilde{\Omega}$ be such that $|y_1-y_2|=r<\dist(y_1,\partial\widetilde{\Omega})$. Consider the ring condenser $(y_1, B(y_1,r))\in \widetilde\Omega$. Then \cite{GResh}
$$
\cp_p(y_1, B(y_1,r))=c_1 r^{n-p}
$$
and 
$$
\cp_q(\varphi^{-1}(y_1), \varphi^{-1}(B(y_1,r)))\geq \cp_q(\varphi^{-1}(y_1), B(\varphi^{-1}(y_1),R))=c_2 R^{n-q},
$$
where $R=\max_{|y_1-y_2|}|\varphi^{-1}(y_1)-\varphi^{-1}(y_2)|$.
Hence, in the case $p=q$, by Theorem~\ref{thm:CapacityDescPP} we obtain
\begin{multline*}
c_2 |\varphi^{-1}(y_1)-\varphi^{-1}(y_2)|^{\frac{n-p}{p}}\leq \cp^{\frac{1}{p}}_p(\varphi^{-1}(y_1), \varphi^{-1}(B(y_1,r)))\\
\leq K_p
\cp^{\frac{1}{p}}_p(y_1, B(y_1,r))=c_1 y_1-y_2|^{\frac{n-p}{p}}.
\end{multline*}
Therefore 
$$
|y_1-y_2|^{\frac{p-n}{p}}\leq c(n,p,q)K_p|\varphi^{-1}(y_1)-\varphi^{-1}(y_2)|^{\frac{p-n}{p}}.
$$

In the case $p>q$, by Theorem~\ref{thm:CapacityDescPQ} we obtain
\begin{multline*}
c_2 |\varphi^{-1}(y_1)-\varphi^{-1}(y_2)|^{\frac{n-q}{q}}\leq \cp^{\frac{1}{q}}_q(\varphi^{-1}(y_1), \varphi^{-1}(B(y_1,r)))\\
\leq \Phi\left((B(y_1,r)\right)^{\frac{p-q}{pq}}
\cp^{\frac{1}{p}}_p(y_1, B(y_1,r))=c_1 y_1-y_2|^{\frac{n-p}{p}}.
\end{multline*}
Therefore 
$$
|y_1-y_2|^{\frac{p-n}{p}}\leq c(n,p,q)\Phi\left((B(y_1,r)\right)^{\frac{p-q}{pq}}|\varphi^{-1}(y_1)-\varphi^{-1}(y_2)|^{\frac{q-n}{q}}.
$$
\end{proof}

\begin{cor}
Let $\varphi:\Omega\to\widetilde{\Omega}$ be a weak $(p,q)$-quasiconformal mapping, $n<q\leq p<\infty$. Then $\varphi$ belongs to the H\"older space $H^{\gamma}(\Omega)$ with $0<\gamma=\frac{q-n}{p-n}\frac{p}{q}\leq 1$.
\end{cor}

In the case $n-1<q\leq p<n$ we have

\begin{thm}
Let $\varphi:\Omega\to\widetilde{\Omega}$ be a weak $(p,q)$-quasiconformal mapping, $n-1<q\leq p< n$. Then for any points $x_1,x_2\in\Omega$, such that $|x_1-x_2|=r<\dist(x_1,\partial\Omega)$ the following inequality
$$
|x_1-x_2|^{\frac{q'-n}{q'}}\leq c(n,p',q')\Phi\left((B(x_1,r)\right)^{\frac{q'-p'}{q'p'}}|\varphi(x_1)-\varphi(x_2)|^{\frac{p'-n}{p'}},
$$
where $p'=p/(p-n+1)$ and $q'=q/(q-n+1)$, holds.
\end{thm}

\begin{proof}
Let Let $\varphi:\Omega\to\widetilde{\Omega}$ be a weak $(p,q)$-quasiconformal mapping, $n-1<q\leq p< n$. Then by the composition duality theorem (Theorem~\ref{CompDual}) we have that the inverse mapping $\varphi^{-1}:\widetilde{\Omega} \to\Omega$ be a weak $(q',p')$-quasiconformal mapping, $n<p'\leq q'< \infty$. Hence by Theorem~\ref{Hol} for any points $x_1,x_2\in\Omega$, such that $|x_1-x_2|=r<\dist(x_1,\partial\Omega)$ the following inequality
$$
|x_1-x_2|^{\frac{q'-n}{q'}}\leq c(n,p',q')\Phi\left((B(x_1,r)\right)^{\frac{q'-p'}{q'p'}}|\varphi(x_1)-\varphi(x_2)|^{\frac{p'-n}{p'}},
$$
where $p'=p/(p-n+1)$ and $q'=q/(q-n+1)$, holds.
\end{proof}

\begin{cor}
Let $\varphi:\Omega\to\widetilde{\Omega}$ be a weak $(p,q)$-quasiconformal mapping, $n-1<q\leq p< n$. Then the inverse mapping $\varphi^{-1}$ belongs to the H\"older space $H^{\gamma}(\widetilde{\Omega})$ with $0<\gamma=\frac{p'-n}{q'-n}\frac{q'}{p'}\leq 1$.
\end{cor}

\vskip 0.2cm

The measure distortion properties of weak $(p,q)$-quasiconformal mappings are based on the "inverse" dilatation function \cite{VU02,VU04} and we do not present them in this review. 

\vskip 0.2cm

\begin{thebibliography}{References}

\bibitem{Bu98} V.~I.~Burenkov, Sobolev Spaces on Domains, Teubner-Texter zur Mathematik, Stuttgart, (1998).

\bibitem{Ch54} G.~Choquet, Theory of capacities, Ann. Inst. Fourier, 155 (1954), 131--295.

\bibitem{CHM}  M.~Cs\"ornyei, S.~Hencl, J.~Mal\'y, Homeomorphisms in the Sobolev space $W^{1,n-1}$,
J. Reine Angew. Math., 644 (2010), 221--235.

\bibitem{DSh58} N.~Dunford and J.~T.~Schwartz, Linear Operators, Part 1: General Theory. Interscience, New York, London, (1958).

\bibitem{Fe69} H.~Federer, Geometric measure theory, Sp\-rin\-ger Verlag, Berlin, (1969).

\bibitem{GPS18} M.Carozza, F.Gianetti, A. Passarelli di Napoli, C.Sbordone, R.Chiattarella, Bi-Sobolev mappings and $K_p$-distortion in the plane. J. Math. Anal. Appl. 457 (2018) 1232-1240.

\bibitem{Ge69} F.~W.~Gehring, Lipschitz mappings and the $p$-capacity of rings in $n$-space, Advances in the theory of Riemann surfaces (Proc. Conf., Stony Brook, N. Y., 1969), 175--193. Ann. of Math. Studies, No. 66. Princeton Univ. Press, Princeton, N. J. (1971)

\bibitem{GR84} V.~Gol'dshtein, A.~S.~Romanov, On mappings preserving the Sobolev spaces. Siberian Math. J., 25, (1984), 55--61.

\bibitem{GGu94} V.~Gol'dshtein, L.~Gurov, Applications of change of variables operators for exact embedding theorems, Integral Equations Operator Theory 19 (1994), 1--24.

\bibitem{GGR95} V.~Gol'dshtein, L.~Gurov, A.~Romanov, Homeomorphisms that induce monomorphisms of Sobolev spaces, 
Israel J. Math., 91 (1995), 31--60.

\bibitem{GPU18_3} V.~Gol'dshtein, V.~Pchelintsev, A.~Ukhlov, On the First Eigenvalue of the Degenerate p-Laplace Operator in Non-convex Domains. Integral Equations Operator Theory 90 (2018), 90:43.

\bibitem{GS82} V.~M.~Gol'dshtein, V.~N.~Sitnikov, Continuation of functions of the class $W^1_p$ across H\"older boundaries, Imbedding theorems and their applications, 31-43, Trudy Sem. S. L. Soboleva, No. 1, (1982), Akad. Nauk SSSR
Sibirsk. Otdel., Inst. Mat., Novosibirsk, (1982).

\bibitem{GU09} V.~Gol'dshtein, A.~Ukhlov, Weighted Sobolev spaces and embedding theorems, Trans. Amer. Math. Soc., 361 (2009), 3829--3850.

\bibitem{GU10-1} V.~Gol'dshtein, A.~Ukhlov, Sobolev homeomorphisms and composition operators. Around the research of Vladimir Maz'ya. I, Int. Math. Ser. (N. Y.), 11 (2010), 207--220.

\bibitem{GU10-2} V.~Gol'dshtein, A.~Ukhlov, About homeomorphisms that induce composition operators on Sobolev spaces, 
Complex Var. Elliptic Equ., 55 (2010), 833--845.

\bibitem{GU16} V.~Gol'dshtein, A.~Ukhlov, On the first Eigenvalues of Free Vibrating Membranes in Conformal Regular Domains. 
Arch. Rational Mech. Anal., 221 (2016), no. 2, 893--915.

\bibitem{GU17} V.~Gol'dshtein, A.~Ukhlov, The spectral estimates for the Neumann-Laplace operator in space domains. 
Adv. in Math., 315 (2017), 166--193.

\bibitem{GResh} V. M. Gol'dshtein, Yu. G. Reshetnyak, Quasiconformal mappings and Sobolev spaces, Dordrecht, Boston, London: Kluwer Academic Publishers, (1990).

\bibitem{Ha93} P.~Hajlasz, Change of variable formula under minimal assumptions, Colloq. Math. 64, (1993), 93--101.

\bibitem{HM72} V.~Havin, V.~Maz'ya, Nonlinear potential theory, Russian Math. Surveys, 27 (1972), 71--148.

\bibitem{HKM} J.~Heinonen, T.~Kilpelinen, O.~Martio, Nonlinear Potential Theory of Degenerate Elliptic Equations. Clarendon Press. Oxford, New York, Tokio (1993).

\bibitem{HK06} S.~Hencl, P.~Koskela, Regularity of the inverse of a planar Sobolev homeomorphism,
Arch. Ration. Mech. Anal., 180 (2006), 75--95.

\bibitem{HKM06} S.~Hencl, P.~Koskela, Y.~Maly, Regularity of the inverse of a Sobolev homeomorphism in space,
Proc. Roy. Soc. Edinburgh Sect. A., 136A (2006), 1267--1285.

\bibitem{HK12} S.~Hencl, L.Kleprik, Composition of q-quasiconformal mappingsand functions of Orlicz-Sobolev spaces. Illinois J. Math. 56(3) (2012) 231-237.

\bibitem{IOZ20}T.~Iwaniec, J.~Onninen, Z.~Zhu, Deformations of bi-conformal energy and a new characterization of quasiconformality, J. Math.Anal.Appl., 457(2018), 1232-1246.

\bibitem{K12} L. Kleprik. Mappings of finite signed distortion . Sobolev spaces and composition of mappings. J. Math. Anal. Appl. 386 (2012)870-881

\bibitem{KRSS14} D.~ A.~Kovtonyuk, V.~I.~Ryazanov, R.~R.~Salimov, E.~A.~Sevost'yanov, Toward the theory of Orlicz-Sobolev classes,
St. Petersburg Math. Journal, 23 (2014), 929--963.

\bibitem{KRSS14M} D.~ A.~Kovtonyuk, V.~I.~Ryazanov, R.~R.~Salimov, E.~A.~Sevost'yanov, Boundary Behavior of Orlicz-Sobolev Classes, Math. Notes, 95 (2014), 509--519.

\bibitem{Kl12} L.~Kleprlik, Mappings of finite signed distortion: Sobolev spaces and composition of mappings, J. Math. Anal. Appl., 386 (2012) 870--881.

\bibitem{K86} V.~I.~Kruglikov, Capacities of condensors and quasiconformal in the mean mappings in space. Mat. Sb. (N.S.) 130 (172) (1986), 185--206.

\bibitem{L71} L.~G.~Lewis, Quasiconformal mappings and Royden algebras in space, Trans. Amer. Math. Soc., 158 (1971), 481--492.

\bibitem{MRV69} O.~Martio, S.~Rickman, J.~V\"ais\"al\"a, Definitions for quasiregular mappings, Ann. Acad. Sei. Fenn., Ser. AI, 448 (1969),1--40.

\bibitem{MRSY} O.~Martio, V.~Ryazanov, U.~Srebro, E.~Yakubov, Moduli in modern mapping theory, Springer Monographs in Mathematics. Springer, New York, 2009.

\bibitem{M} V.~Maz'ya, Sobolev spaces: with applications to elliptic partial differential equations, Springer: Berlin/Heidelberg, (2010).

\bibitem{MS86} V.~Maz'ya, T.~O.~Shaposhnikova, Multipliers in Spaces of Differentiable Functions, 
Leningrad Univ. Press., 1986.

\bibitem{M90} I.~G.~Markina, A change of variable preserving the differential properties of functions. Siberian Math. J., 31, (1990), 73--84.

\bibitem{MRSY01} O. Martio, V. Ryazanov, U. Srebro, E. Yakubov, To the theory of $Q$-homeomorphisms, Dokl. Akad. Nauk Rossii, 381 (2001), 20-22.

\bibitem{MU21} A.~Menovschikov, A.~Ukhlov, Composition operators on Sobolev spaces, $Q$-mappings and weighted Sobolev inequalities,  	arXiv:2110.09261.

\bibitem{N60} M.~Nakai, Algebraic criterion on quasiconformal equivalence of Riemann surfaces, Nagoya Math. J., 16 (1960), 157--184.

\bibitem{OS65} I.~S.~Ovchinnikov, G.~D.~Suvorov, Transformations of the Dirichlet integral and space mappings, 
Siberian Math. J., 6 (1965), 1292--1314.

\bibitem{PW} L.~E.~Payne, H.~F.~Weinberger, An optimal Poincar\'e inequality for convex domains, 
Arch. Rat. Mech. Anal., 5 (1960), 286--292.

\bibitem{P93} Yu.~A.~Peshkichev, Inverse mappings for homeomorphisms of the class $BL$, Mathematical Notes. 53 (1993), 98--101.

\bibitem{PS51} G.~P\'olya, G. Szeg\"o, Isoperimetric Inequalities in Mathematical Physics, Princeton University Press, 1951.

\bibitem{RR55} T.~Rado, P.~V.~Reichelderfer, Continuous Transformations in Analysis. Springer-Verlag, Berlin (1955).

\bibitem{Sa15} R.~R.~Salimov, Lower Q-homeomorphisms with respect to the p-modulus, Ukr. Mat. Visn. 12 (2015),484--510. 

\bibitem{S16} E.~Sevost'yanov, On open and discrete mappings with a modulus condition, Ann. Acad. Sci. Fenn. Math. 41 (2016), 41--50.

\bibitem{S41} S.~L.~Sobolev, Some groups of transformations of an n-dimensional space, Dokl. Akad. Nauk SSSR, 32 (1941), 380--382.

\bibitem{U93} A.~Ukhlov, On mappings, which induce embeddings of Sobolev spaces, Siberian Math. J., 34 (1993), 185--192.

\bibitem{U04} A.~Ukhlov, Differential and geometrical properties of Sobolev mappings, Matem. Notes., 75 (2004), 291--294. 

\bibitem{UV10} A.~Ukhlov, S.~K.~Vodop'yanov, Mappings with bounded $(P,Q)$-distortion on Carnot groups, Bull. Sci. Math., 134 (2010), 605--634.

\bibitem{Va} Ju.~V\"ais\"al\"a, Lectires on $n$-dimensional quasiconformal mappings, Springer Verlag, (1955).

\bibitem{V88} S.~K.~Vodop'yanov, Taylor Formula and Function Spaces, Novosibirsk, Novosibirsk Univ. Press., 1988.

\bibitem{V89} S.~K.~Vodop'yanov, Mappings of homogeneous groups and embeddings of function spaces, Siberian Math. J., 30 (1989), 25--41.

\bibitem{V90} S.~K.~Vodop'yanov, $L_p$-theory of potential and quasiconformal mappings on homogeneous groups,
Sovrem. Probl. Geom. Analiz., 263 (1990), 45--89. 

\bibitem{V05} S.~K.~Vodop'yanov, Composition operators on Sobolev spaces, Complex analysis and dynamical systems II (Nahariya, Israel, 2003), Contemp. Math., 382, Amer. Math. Soc., Providence, RI, (2005), 401--415.

\bibitem{V12} S.~K.~Vodop'yanov, Regularity of mappings inverse to Sobolev mappings, Sbornik: Mathematics, 203 (2012), 1383--1410.

\bibitem{VG75} S.~K.~Vodop'yanov, V.~M.~Gol'dstein, Lattice isomorphisms of the spaces $W^1_n$ and quasiconformal mappings, Siberian Math. J., 16 (1975), 224--246.

%\bibitem{VG75_2} S. K. Vodop'yanov, V. M. Gol'dstein, A criterion for the removability of sets for $W^{1}_{p}$ spaces of quasiconformal and %quasi-isometric mappings, Dokl. Akad. Nauk SSSR, 220 (1975), 769--771.

\bibitem{VG76}S.~K.~Vodop'yanov, V.~M.~Gol'dshtein, Functional characteristics of quasiisometric mappings. Sib. Mat. Zh. 17,(1976),
768--773.

\bibitem{VG77} S. K. Vodop'yanov, V. M. Gol'dstein, Criteria for the removability of sets for $L^1_p$ spaces of quasiconformal and quasi-isomorphic mappings, Siberian Math. J., 18 (1977), 35--50.

\bibitem{VGR79} S.~K.~Vodop'yanov, V.~M.~Gol'dshtein, Yu.~G.~Reshetnyak, On geometric properties of functions with generalized first derivatives, Uspekhi Mat. Nauk, 34 (1979), 17--65. 

\bibitem{VU98} S.~K.~Vodop'yanov, A.~D.~Ukhlov, Sobolev spaces and $(P,Q)$-quasiconformal mappings of Carnot groups.
Siberian Math. J., 39 (1998), 665--682.

\bibitem{VU02} S.~K.~Vodop'yanov, A.~D.~Ukhlov, Superposition operators in Sobolev spaces. Russian Mathematics (Iz. VUZ)., 46 (2002), 11--33. 

\bibitem{VU04} S.~K.~Vodop'yanov, A.~D.~Ukhlov, Set functions and their applications in the theory of Lebesgue and Sobolev spaces, 
Siberian Adv. in Math, 14 (2004), 78--125.

\bibitem{VU05} S.~K.~Vodop'yanov, A.~D.~Ukhlov, Set functions and their applications in the theory of Lebesgue and Sobolev spaces, 
Siberian Adv. in Math, 15 (2005), 91--125.

\bibitem{Yo} K.~Yosida, Functional Analysis, Springer Verlag, Berlin, (1980).

\bibitem{Z69} Ziemer W.~P.: Change of variables for absolutely continuous functions. Duke Math. J. 36, 171--178  (1969)

\bibitem{Zi} W.~P.~Ziemer, Weakly Differentiable Functions, Springer-Verlag, New York, (1989).

\end{thebibliography}

\noindent
Vladimir Gol'dshtein  \,  \hskip 3.2cm Alexander Ukhlov

\noindent
Department of Mathematics   \hskip 2.25cm Department of Mathematics

\noindent
Ben-Gurion University of the Negev  \hskip 1.05cm Ben-Gurion University of the Negev

\noindent
P.O.Box 653, Beer Sheva, 84105, Israel  \hskip 0.7cm P.O.Box 653, Beer Sheva, 84105, Israel

\noindent
E-mail: vladimir@bgu.ac.il  \hskip 2.5cm E-mail: ukhlov@math.bgu.ac.il
\end{document}



